\chapter{Representations of the coupling matrix}\label{Trep}
In this chapter we recommence the practice of writing the material 
properties as the sum of a reference and residual component, i.e.
$(\rho, \alpha, \beta)=(\rf{\rho}+{\mit \Delta}\rho,\rf{\alpha}+{\mit 
\Delta}\alpha,\rf{\beta}+{\mit \Delta}\beta)$ where the  uniform host 
material has the reference   properties $(\rf{\rho},\rf{\alpha},\rf{\beta})$.
 



\section{Overview}
In chapter \ref{Grep}, we developed various  plane wave representations of the 
displacement/ traction elastic wave Green's tensor ${\rf{\bmi G}}$, 
appropriate for a lattice of source points. We were purposely vague about the nature of the 
source ${\bmi f}_j$ at each lattice point $1 \le j \le L N_x$ other than 
to suppose that it was a $6-$vector. 
In this chapter we will see that some of the components of ${\bmi f}_j$ 
have $\partial_x$ operators. We always transfer these partial 
differentiations to the source and receiver coordinates of the relevant 
Green's tensors, using integration by parts. Consequently,  it is convenient to augment 
the definitions of the Green's tensors to carry these partially 
differentiated components as well as the usual displacement and traction 
fields. The Green's tensors are then expressed through $8 \times 8 $ matrices,
rather than the $6 \times 6$ arrays that we have employed in  chapter 
\ref{Grep}. The advantage of this representation is that the scattering 
induced by heterogeneity superimposed on the reference medium may be 
modelled as a {\em multiplicative operator\/} rather than the 
differential operator (\ref{dAdef}). This is the essence of the {\em 
pseudo-spectral approximation\/} (Orszag, 1972), where a {\em spectral 
convolution\/} is replaced by a {\em spatial multiplication\/}. 

We recall that the convolution of two $N-$sequences, is a $2 N-$sequence.
The ``approximation" 
in the term ``pseudo-spectral approximation" derives from the neglect of 
the aliasing terms that result from sampling the $2N-$spectral 
convolution (of two $N-$spectral sequences) as if it was an  $N-$sequence-- 
the approximation that is commonly employed in what has become known as the pseudo-spectral 
method [e.g. {\em Fornberg\/}~(1987)]. In our method, we are 
oversampling by a factor $\ge 2$ (usually $16$, or $32$). So that the 
aliasing contribution we have just mentioned will not occur. 


In Corollary 1, of \S\ref{secLS} we found the singular 
integro/differential {\em Lippmann-Schwinger type operator 
equation\/}, reproduced here as
\begin{eqnarray} \label{LSTrep}
{{\bmi G}}({\bmi x};{\bmi x}_{S};k_y)&=&
{\rf{\bmi G}}({\bmi x};{\bmi x}_{S};k_y)\nonumber\\
&&+\int_{{\mit \Omega}} {\rf{\bmi G}}({\bmi x};{\bmi x}^{\prime};k_y)
{\mit \Delta} {\bmi A}({\bmi
x}^{\prime},\partial_{{x}^{\prime}};k_y) {\bmi G}({\bmi x}^{\prime},{\bmi x}_{S};k_y) {\rm d}{\bmi 
x}^{\prime},
\quad {\bmi x}\in R^{2}
\end{eqnarray}
for the unknown Green's tensor ${{\bmi G}}({\bmi x};{\bmi 
x}_{S};k_y)$, defined in terms of the {\em known reference medium Green's 
tensor\/} 
${\rf{\bmi G}}$ where the {\em differential coupling operator\/}  ${{\mit 
\Delta} {\bmi A}}$ is 
defined as 
\begin{samepage}
\begin{eqnarray*}
\lefteqn{
{{\mit \Delta} {\bmi A}}({\bmi x};{\partial }_{x},{k}_{y})
=\hspace{26em}(\ref{dAdef})}\\
&&\!\!\!\!\!\!\!\!\!\!\!\!
\left[{\begin{array}{cccccc}
0&- {{\mit \Delta} \gamma} {\partial }_{  x}&  
-i {{\mit \Delta} \gamma} {  k}_{  y}&  {{\mit \Delta} a}&0&0\\0&0&0&0&{{\mit \Delta} b}&0\\0&0&0&0&0&{{\mit \Delta} b}\\ - {{\mit \Delta} \rho} {\omega }^{  
2}&  0&0&0&0&0\\ 0&- {{\mit \Delta} \rho} {\omega 
}^{  2}  -{ \partial }_{x} {{\mit \Delta} \zeta} {\partial }_{  
x}  + {{\mit \Delta} \mu} {  k}_{  y}^{  2}&  
-i{k}_{y}{ \partial }_{x} {{\mit \Delta} \chi}   -i{k}_{y} {{\mit \Delta} \mu} {\partial 
}_{  x}&  -{ \partial }_{x} {{\mit \Delta} \gamma} &  0&0\\ 
0&-i{k}_{y} {{\mit \Delta} \chi} {\partial }_{  x}  -i{k}_{y}{ \partial 
}_{x} {{\mit \Delta} \mu} &  - {{\mit \Delta} \rho} {\omega }^{  2}  + {{\mit \Delta} \zeta }
{  k}_{  y}^{  2}  -{ \partial }_{x} {{\mit \Delta} \mu} 
{\partial }_{  x}&  -i{k}_{y} {{\mit \Delta} \gamma} &  0&0\end{array}}\right]
\end{eqnarray*}
\end{samepage}
with the material contrast parameters in (\ref{dAdef}) defined in Table \ref{contrastTable}.
 The integral $\int_{{\mit \Omega}} \cdot$ in Corollary 1
was defined over the spatial support ${\mit \Omega}$ of the inclusion.

In this chapter, we will replace the singular 
integro/differential {\em Lippmann- Schwinger type operator 
equation\/} of chapter \ref{MSrep}  by a nonsingular, compact integral operator equation, that 
is much simpler to discretize and solve. To effect this transformation we
employ a {\em Galerkin\/} type projection 
of the range space of (\ref{LSTrep}) onto the span of 
\[
\int_{{\mit \Omega}}{\rm d}{\bmi x}\, {\rf{\bmi G}}({\bmi x}_{R};{\bmi x};k_y)
{\mit \Delta} {\bmi A}({\bmi
x},\partial_{x};k_y) \cdot
\]
so that (\ref{LSTrep}) is replaced by
\begin{eqnarray} \label{LSTrep1}
\lefteqn{\int_{{\mit \Omega}} {\rm d}{\bmi x}\,{\rf{\bmi G}}({\bmi x}_{R};{\bmi x};k_y)
{\mit \Delta} {\bmi A}({\bmi
x},\partial_{x};k_y){{\bmi G}}({\bmi x};{\bmi x}_{S};k_y)=}\nonumber\\
&&
\int_{{\mit \Omega}} {\rm d}{\bmi x}\,{\rf{\bmi G}}({\bmi x}_{R};{\bmi x};k_y)
{\mit \Delta} {\bmi A}({\bmi
x},\partial_{x};k_y){\rf{\bmi G}}({\bmi x};{\bmi x}_{S};k_y)\\
&&+\int_{{\mit \Omega}} {\rm d}{\bmi x}\, {\rf{\bmi G}}({\bmi x}_{R};{\bmi x};k_y)
{\mit \Delta} {\bmi A}({\bmi
x},\partial_{x};k_y)\int_{{\mit \Omega}}{\rm d}{\bmi x}^{\prime}\, 
{\rf{\bmi G}}({\bmi x};{\bmi x}^{\prime};k_y)
{\mit \Delta} {\bmi A}({\bmi
x}^{\prime},\partial_{x^{\prime}};k_y) {\bmi G}({\bmi x}^{\prime},{\bmi 
x}_{S};k_y) .
\nonumber
\end{eqnarray}

Let us denote the reference ${\rf{\bmi G}}$ and actual ${{\bmi G}}$ 
displacement/traction Green's tensors in terms of their $3 \times 3$ 
displacement and traction partitions  as
\[
{\rf{\bmi G}}=\left[{\begin{array}{cc}
\rf{\bmi P}_{uu}&\rf{\bmi P}_{ut}\\ 
\rf{\bmi P}_{tu}&\rf{\bmi P}_{tt}
\end{array}}\right]{\bmi J}_{6}, \quad \mbox{and} \quad
{{\bmi G}}=\left[{\begin{array}{cc}
{\bmi P}_{uu}&{\bmi P}_{ut}\\ 
{\bmi P}_{tu}&{\bmi P}_{tt}
\end{array}}\right]{\bmi J}_{6}
\] 
respectively. 
Next we employ integration by parts in (\ref{LSTrep1}) to transfer the 
$\partial_{x}$ operators  in ${\mit \Delta} {\bmi A}$ onto the field 
coordinate of the unknown Green's tensor  ${{\bmi G}}$ to give the augmented 
operator ${{\bmi G}_{r}}$, and onto both the 
source and field coordinate of the reference  Green's tensor  ${\rf{\bmi 
G}}$, to give ${\rf{\bmi P}_{r,s}}$ where
the augmented operators may be written as:
\[{{\bmi G}_{r}}=\left[{\begin{array}{cc}{\bmi P}_{uu}&{\bmi P}_{ut}\\ 
{\bmi P}_{tu}&{\bmi P}_{tt}\\
{\bmi P}_{du}&{\bmi P}_{dt}\end{array}}\right]{\bmi J}_{6}
, \quad
{\rf{\bmi P}_{s}}=\left[{\begin{array}{ccc}
\rf{\bmi P}_{uu}&\rf{\bmi P}_{ut}&\rf{\bmi P}_{ud}\\ 
\rf{\bmi P}_{tu}&\rf{\bmi P}_{tt}&\rf{\bmi P}_{td}
\end{array}}\right]
, \quad
{\rf{\bmi P}_{r,s}}=\left[{\begin{array}{ccc}
\rf{\bmi P}_{uu}&\rf{\bmi P}_{ut}&\rf{\bmi P}_{ud}\\ 
\rf{\bmi P}_{tu}&\rf{\bmi P}_{tt}&\rf{\bmi P}_{td}\\
\rf{\bmi P}_{du}&\rf{\bmi P}_{dt}&\rf{\bmi P}_{dd}
\end{array}}\right]
\]
 These integration by parts 
operations may be summarized in the somewhat stylized notation 
\begin{samepage}
\begin{eqnarray} \label{LSTrep2}
\lefteqn{\int_{{\mit \Omega}} {\rm d}{\bmi x}\,{\rf{\bmi P}}_{s}({\bmi x}_{R};{\bmi x};k_y)
{\mit \Delta} {\bmi C}({\bmi
x};k_y){{\bmi G}}_{r}({\bmi x};{\bmi x}_{S};k_y)=}\nonumber\\
&&
\int_{{\mit \Omega}} {\rm d}{\bmi x}\,{\rf{\bmi P}}_{s}({\bmi x}_{R};{\bmi x};k_y)
{\mit \Delta} {\bmi C}({\bmi
x};k_y){\rf{\bmi G}}_{r}({\bmi x};{\bmi x}_{S};k_y)\\
&&+\int_{{\mit \Omega}} {\rm d}{\bmi x}\, {\rf{\bmi P}}_{s}({\bmi x}_{R};{\bmi x};k_y)
{\mit \Delta} {\bmi C}({\bmi
x};k_y)\int_{{\mit \Omega}}{\rm d}{\bmi x}^{\prime}\, 
{\rf{\bmi P}}_{r,s}({\bmi x};{\bmi x}^{\prime};k_y)
{\mit \Delta} {\bmi C}({\bmi
x}^{\prime};k_y) {\bmi G}_{r}({\bmi x}^{\prime},{\bmi 
x}_{S};k_y) .
\nonumber
\end{eqnarray}
\end{samepage}
where ${\mit \Delta} {\bmi C}$ is an $8 \times 8$ multiplicative operator 
defined as
\begin{samepage}
\begin{eqnarray*}
\lefteqn{
{\mit \Delta}{\bmi C}({\bmi x};k_y)=}\\
&&
\left[{\begin{array}{cccccccc}
\displaystyle - {\mit \Delta}\rho {\omega }^{ 2}& \displaystyle 0& 
\displaystyle 0& \displaystyle 0& \displaystyle 0& \displaystyle 0& 
\displaystyle 0& \displaystyle 0\\ \displaystyle 0& \displaystyle {\mit \Delta}\mu { 
k}_{ y}^{ 2} - {\mit \Delta}\rho {\omega }^{ 2}& \displaystyle 0& \displaystyle 0& 
\displaystyle 0& \displaystyle 0& \displaystyle 0& \displaystyle -i{k}_{y} 
{\mit \Delta}\mu \\ \displaystyle 0& \displaystyle 0& \displaystyle {\mit \Delta}\zeta { k}_{ y}^{ 2} 
- {\mit \Delta}\rho {\omega }^{ 2}& \displaystyle -i{k}_{y} {\mit \Delta}\gamma & \displaystyle 0& 
\displaystyle 0& \displaystyle -i{k}_{y} \chi & \displaystyle 0\\ 
\displaystyle 0& \displaystyle 0& \displaystyle i {\mit \Delta}\gamma { k}_{ y}& 
\displaystyle -{\mit \Delta}a& \displaystyle 0& \displaystyle 0& \displaystyle {\mit \Delta}\gamma & 
\displaystyle 0\\ \displaystyle 0& \displaystyle 0& \displaystyle 0& 
\displaystyle 0& \displaystyle -{\mit \Delta}b& \displaystyle 0& \displaystyle 0& 
\displaystyle 0\\ \displaystyle 0& \displaystyle 0& \displaystyle 0& 
\displaystyle 0& \displaystyle 0& \displaystyle -{\mit \Delta}b& \displaystyle 0& 
\displaystyle 0\\ \displaystyle 0& \displaystyle 0& \displaystyle i{k}_{y} 
{\mit \Delta}\chi & \displaystyle {\mit \Delta}\gamma & \displaystyle 0& \displaystyle 0& 
\displaystyle {\mit \Delta}\zeta & \displaystyle 0\\ \displaystyle 0& \displaystyle 
i{k}_{y} {\mit \Delta}\mu & \displaystyle 0& \displaystyle 0& \displaystyle 0& 
\displaystyle 0& \displaystyle 0& \displaystyle {\mit \Delta}\mu 
\end{array}}\right]\nonumber 
\end{eqnarray*}
\end{samepage}

However, (\ref{LSTrep2}) is equivalent to 
\begin{samepage}
\begin{eqnarray} \label{LSTrep30}
\lefteqn{\int_{{\mit \Omega}} {\rm d}{\bmi x}\,{\rf{\bmi P}}_{s}({\bmi x}_{R};{\bmi x};k_y)
{\mit \Delta} {\bmi C}({\bmi
x};k_y)\left[{{\bmi G}}_{r}({\bmi x};{\bmi x}_{S};k_y)-{\rf{\bmi G}}_{r}({\bmi x};{\bmi x}_{S};k_y)
\right.}\nonumber\\
&&\qquad \qquad \qquad \qquad \qquad
\left.-\int_{{\mit \Omega}}{\rm d}{\bmi x}^{\prime}\, 
{\rf{\bmi P}}_{r,s}({\bmi x};{\bmi x}^{\prime};k_y)
{\mit \Delta} {\bmi C}({\bmi
x}^{\prime};k_y) {\bmi G}_{r}({\bmi x}^{\prime},{\bmi 
x}_{S};k_y)\right]= {\bf 0} .
\nonumber
\end{eqnarray}
\end{samepage}
which is the {\em weak\/} statement of the {\em orthogonality\/} of the residual
\[{\bmi R}\equiv
{{\bmi G}}_{r}({\bmi x};{\bmi x}_{S};k_y)-{\rf{\bmi G}}_{r}({\bmi x};{\bmi x}_{S};k_y)
-\int_{{\mit \Omega}}{\rm d}{\bmi x}^{\prime}\, 
{\rf{\bmi P}}_{r,s}({\bmi x};{\bmi x}^{\prime};k_y)
{\mit \Delta} {\bmi C}({\bmi
x}^{\prime};k_y) {\bmi G}_{r}({\bmi x}^{\prime},{\bmi 
x}_{S};k_y)
\] 
 to the {\em test\/} space spanned by
\[
\int_{{\mit \Omega}} {\rm d}{\bmi x}\,{\rf{\bmi P}}_{s}({\bmi x}_{R};{\bmi x};k_y)
{\mit \Delta} {\bmi C}({\bmi
x};k_y) \left[\cdot\right], \qquad \mbox{for arbitrary }\quad {\bmi x}_{R}
\]

In this chapter, we will solve for ${{\bmi G}}_{r}$ from the {\em 
augmented equation\/} 
\begin{equation} \label{LSTrep3}
{{\bmi G}}_{r}({\bmi x};{\bmi x}_{S};k_y)={\rf{\bmi G}}_{r}({\bmi x};{\bmi x}_{S};k_y)
+\int_{{\mit \Omega}}{\rm d}{\bmi x}^{\prime}\, 
{\rf{\bmi P}}_{r,s}({\bmi x};{\bmi x}^{\prime};k_y)
{\mit \Delta} {\bmi C}({\bmi
x}^{\prime};k_y) {\bmi G}_{r}({\bmi x}^{\prime},{\bmi 
x}_{S};k_y)
\end{equation} 
instead of ${{\bmi G}}$ from (\ref{LSTrep1}). The advantage of 
(\ref{LSTrep3}), is that it is an integral equation,  while
 the original formulation (\ref{LSTrep1}) was an integro/differential 
 equation. 

However, (\ref{LSTrep3}) is {\em strongly singular\/} because
\[
{\rf{\bmi P}}_{r,s}({\bmi x};{\bmi x}^{\prime};k_y)\approx 
O\left(\| {\bmi x}-{\bmi x}^{\prime} \|^{-2} \right), \quad \mbox{in the 
limit} \quad \| {\bmi x}-{\bmi x}^{\prime} \|\rightarrow 0
\]
We can {\em subtract \/} the singular part of (\ref{LSTrep3}) and invert 
it analytically. The resulting {\em renormalized\/} equation
\begin{eqnarray} \label{LSTrep4}
\lefteqn{{{\bmi G}}_{r}({\bmi x};{\bmi x}_{S};k_y)={\rf{\bmi 
G}}_{r}({\bmi x};{\bmi x}_{S};k_y)}\nonumber\\
&&
+\int_{{\mit \Omega}}{\rm d}{\bmi x}^{\prime}\, 
{\rf{\bmi P}}_{r,s}({\bmi x};{\bmi x}^{\prime};k_y)
{\mit \Delta} {\bmi C}_{eff}({\bmi
x}^{\prime};k_y) {\bmi G}_{r}({\bmi x}^{\prime},{\bmi 
x}_{S};k_y), \quad {\bmi x}\not={\bmi x}^{\prime}
\end{eqnarray} 
is both {\em nonsingular\/} and {\em compact\/}, where 
we have introduced the {\em effective coupling matrix\/} ${\mit \Delta} 
{\bmi C}_{eff}$ defined later in (\ref{dCeff}).

%%%%%%%%%%%%%%%%%%%%%%%%%%%%%%%%%%%%%%%%%%%%%%%%%%%%%%%%%%%%%%%%%%%%%%%%%%%%%%%%%%%%%%%%%%%
%%%%%%%%%%%%%%%%%%%%%%%%%%%%%%%%%%%%%%%%%%%%%%%%%%%%%%%%%%%%%%%%%%%%%%%%%%%%%%%%%%%%%%%%%%%
%%%%%%%%%%%%%%%%%%%%%%%%%%%%%%%%%%%%%%%%%%%%%%%%%%%%%%%%%%%%%%%%%%%%%%%%%%%%%%%%%%%%%%%%%%%
%%%%%%%%%%%%%%%%%%%%%%%%%%%%%%%%%%%%%%%%%%%%%%%%%%%%%%%%%%%%%%%%%%%%%%%%%%%%%%%%%%%%%%%%%%%
\section{The two centred, harmonic expansion of the reference Green's tensor}
In this section, we derive a closed form expression for the 
reference medium Green's tensor $\rf{\bmi G}$, in a {\em basis of cylindrical harmonics\/}. 
The expansion is two centred, i.e. it is constructed in two separate coordinate 
systems, one coordinate system for the {\em source\/} point
${\bmi  x}^{\prime}$, and one coordinate system for the {\em field\/} point  ${\bmi  x}$.
We will utilize $\rf{\bmi G}$ to investigate the near field 
singularity of ${\rf{\bmi P}}_{r,s}$ from (\ref{LSTrep3}).

We recall from \S\ref{secxsplit}, that
$\rf{\bmi G}$ may be written as 
\begin{samepage} 
\begin{eqnarray}\label{c4grepc}
\rf{\bmi  G}({\bmi  x};{{\bmi  x}}^{\prime};k_y)& =&\sum\limits_{c\in \{P,S,H\}}
\frac{1}{4{ \pi }^{2}\rho \omega^2}
\int_{{R}^{2}} {\rm d}{\bmi k}\, {\bmi L}_{c} \otimes {\bmi  L}_{c}^{\prime}
\frac{
\exp \left[{i {\bmi k}  \cdot  ({\bmi x}-{\bmi x}^{\prime})}\right]  
}{ \hat{K}^{2}_{ c } - {k}^{2}_{z} - {k}^{2}_{x}} {\bmi J}_{6}\nonumber\\
&=&
\rf{\bmi  G}_{P}({\bmi  x};{{\bmi  x}}^{\prime};k_y)+
\rf{\bmi  G}_{S}({\bmi  x};{{\bmi  x}}^{\prime};k_y)+
\rf{\bmi  G}_{H}({\bmi  x};{{\bmi  x}}^{\prime};k_y)
\end{eqnarray} 
\end{samepage}
where  $\otimes$ is the {\em tensor product\/} and 
 the partial differential 
operators
\begin{eqnarray} \label{c4Labdef}
\begin{array}{lll}
{\bmi L}_{P}=&{\bmi L}_{S}=&{\bmi L}_{H}=\\
\left[\begin{array}{c}
\partial_{z} \\
\partial_{x} \\ 
\partial_{y} \\ 
2 \rf{\mu} \partial_{zz} - {\lambda \omega^{2}\over \rf{\alpha}\!{}^2}\\
    2 \rf{\mu} \partial_{xz} \\ 
   2 \rf{\mu} \partial_{yz}
\end{array}\right], \quad& 
\frac{\omega }{\rf{\beta} \hat{K}_{S}}\,
\left[\begin{array}{c}
\partial_{x} \\-\partial_{z} \\ 0 \\
2 \rf{\mu} \partial_{xz}\\  
   \rf{\mu} (-\partial_{zz} + \partial_{xx}) \\ 
   \rf{\mu} \partial_{xy} 
\end{array}\right], \quad&
\frac{1 }{\hat{K}_{H}}\,
\left[\begin{array}{c}
   \partial_{yz} \\
\partial_{xy} \\ 
{\omega^{2}\over \rf{\beta}\!{}^2} + \partial_{yy} \\ 
 2 \rf{\mu} \partial_{yzz}\\
    2 \rf{\mu} \partial_{xyz} \\ 
   \rf{\mu} ({\omega^{2}\over \rf{\beta}\!{}^2}\partial_{z} + 2 \partial_{yyz}) 
\end{array}\right]
\end{array}
\end{eqnarray}
were defined earlier in (\ref{Labdef}).
%%%%%%%%%%%%%%%%%%%%%%%%%%%%%%%%%%%%%%%%%%%%%%%%%%%%%%%%%%%%%%%%%%%%%%%%%%%%%%%%%%
%%%%%%%%%%%%%%%%%%%%%%%%%%%%%%%%%%%%%%%%%%%%%%%%%%%%%%%%%%%%%%%%%%%%%%%%%%%%%%%%%%
%%%%%%%%%%%%%%%%%%%%%%%%%%%%%%%%%%%%%%%%%%%%%%%%%%%%%%%%%%%%%%%%%%%%%%%%%%%%%%%%%%
%%%%%%%%%%%%%%%%%%%%%%%%%%%%%%%%%%%%%%%%%%%%%%%%%%%%%%%%%%%%%%%%%%%%%%%%%%%%%%%%%%

\begin{figure}
\HideDisplacementBoxes
%\ShowDisplacementBoxes
\centerline{\BoxedEPSF{coords.EPSF scaled 1000}}
\caption{A pair of shifted cylindrical coordinate systems, with 
coordinate centres at $\mbox{\mib R}{}_{i}\equiv(Z_i,X_i)$
and $\mbox{\mib R}{}_{j}\equiv(Z_j,X_j)$. The source point
$\mbox{\mib x}{}^{\prime}$
and field 
point $\mbox{\mib x}$ are inside non-overlapping discs. This geometrical 
constraint is often called the {\em muffin tin\/} condition. }
\protect\label{coords}
\end{figure}
%%%%%%%%%%%%%%%%%%%%%%%%%%%%%%%%%%%%%%%%%%%%%%%%%%%%%%%%%%%%%%%%%%%%%%%%%%%%%%%%%%







Assuming that ${\bmi  x}\not={{\bmi  x}}^{\prime}$ in (\ref{c4grepc}), the 
wavenumber integrals are {\em uniformly convergent\/} (due to the choice of complex 
frequency), so that the differential operators may be taken outside the 
respective integrals to get
\begin{eqnarray} \label{grep3}
\rf{\bmi  G}({\bmi  x};{{\bmi  x}}^{\prime};k_y)&=&\sum\limits_{c\in \{P,S,H\}}
 {\bmi L}_{c } \otimes {\bmi  L}_{c}^{\prime}
\left[ \frac{1}{4{ \pi }^{2}\rho \omega^2}
\int_{{R}^{2}} {\rm d}{\bmi k} 
\frac{
\exp \left[{i {\bmi k}  \cdot  ({\bmi x}-{\bmi x}^{\prime})}\right]  
}{ \hat{K}^{2}_{ c } - {k}^{2}_{z} - {k}^{2}_{x}}
\right] {\bmi J}_{6}\nonumber\\
&=&
\rf{\bmi  G}_{P}({\bmi  x};{{\bmi  x}}^{\prime};k_y)+
\rf{\bmi  G}_{S}({\bmi  x};{{\bmi  x}}^{\prime};k_y)+
\rf{\bmi  G}_{H}({\bmi  x};{{\bmi  x}}^{\prime};k_y)
\end{eqnarray}
where $\hat{K}_{c}^2$ was given by (\ref{Kdef}).
Or equivalently, we have
\begin{eqnarray}	\label{Gdef1}
\rf{\bmi  G}({\bmi  x};{{\bmi  x}}^{\prime};k_y)&=& {\frac{1}{\rho 
\omega^2}}\left[{\bmi 
L}_{P} \otimes {\bmi  L}_{P}^{\prime}\left[{ g}_{P } ({\bmi  x}-{{\bmi 
	x}}^{\prime})\right] {\bmi  J}_6 \  
	+\  {\bmi  L}_{S}
	\otimes {\bmi  L}_{S}^{\prime}\left[{ g}_{S } ({\bmi  x}-{{\bmi 
	x}}^{\prime})\right] {\bmi  J}_6\right.\nonumber\\
	&&\left. +\  {\bmi  L}_{H}
	\otimes {\bmi  L}_{H}^{\prime}\left[{ g}_{H } ({\bmi  x}-{{\bmi 
	x}}^{\prime})\right] {\bmi  J}_6 \right], \quad {\bmi  x}\not={{\bmi  x}}^{\prime}
\end{eqnarray}
where the {\em scalar Green's function\/},
\begin{equation} \label{Ig}
{g}_{c }({\bmi  x} -{\bmi  x}^{\prime} ) =  \frac{1}{4{ \pi }^{2}}
\int_{{\Bbb R}^2} \frac{\exp \left[i\ {\bmi  k}
\cdot ({\bmi  x} -{\bmi  x}^{\prime} )\right]}{\hat{K}^{2}_{ c } - {k}^{2}_{z} - {k}^{2}_{x}}\
{ {\rm d}}{\bmi  k},\quad  c\in\{P,S,H\}
\end{equation} 
is the solution of the {\em Helmholtz equation\/}
\begin{equation} \label{HelmhotzEqn}
 {\mit \Delta} {  g}_{c }  ({\bmi  x}-{{\bmi  x}}^{ \prime})+\hat{K}_{ c }^{2}{g}_{ c
}({\bmi  x}-{{\bmi  x}}^{ \prime})\ =\  \delta   ({\bmi  x}-{{\bmi  x}}^{ 
\prime}).
\end{equation}

We have invoked the condition ${\bmi  x}\not={{\bmi  x}}^{\prime}$ in (\ref{c4grepc})
in order to take the differential operators outside the integral signs, 
so that (\ref{Gdef1}) is inapplicable in the limit $\| {\bmi  
x}-{{\bmi  x}}^{\prime} \|\rightarrow 0$. In this limit, ${  g}_{c }$ in 
(\ref{HelmhotzEqn}) is a {\em distribution\/}, so that the $\partial_z, 
\partial_x$ 
operators implicit in some of the components of ${\bmi 
L}_{c} \otimes {\bmi  L}_{c}^{\prime}{  g}_{c }, \ c\in\{P,S,H\}$, will 
generate {\em 
delta functions\/}. We will re-introduce these missing delta function 
terms in \S\ref{secEffScat}.

Introducing the {\em local coordinate transformations\/}({\em vid.\/} {Fig}.\ref{coords})
\[
{\bmi  x}={\bmi  R}_{i}+{\bmi  r}_{i} \quad \mbox{and} \quad 
{\bmi  x}^{\prime}={\bmi  R}_{j}+{\bmi  r}_{j}
\]
for  {\em local coordinate centres\/}  ${\bmi  R}_{i}\equiv(Z_i,X_i)$ and ${\bmi 
R}_{j}\equiv(Z_j,X_j)$, 
(\ref{Ig}) may be
written as
\begin{eqnarray} \label{Ig1} 
\lefteqn{{g}_{c }({\bmi  x} -{\bmi  x}^{\prime} )}\nonumber\\
&&\!\!\!\!\!\!
={\frac{1}{4{ \pi }^{2}}}
\int {\frac{
\exp \left({i\ {\bmi  k}
\cdot ({\bmi  x} -{\bmi  R}_{i})}\right) 
\exp \left({i\ {\bmi  k}
\cdot ({\bmi  R}_{i} -{\bmi  R}_{j})}\right)
\exp \left({-i\ {\bmi  k}
\cdot ({\bmi  x}^{\prime} -{\bmi  R}_{j} )}\right)} 
{\hat{ K}_{c}^{2}-{k}_{z}^{2}-{k}_{x}^{2}}}\ {{\rm d}}{\bmi  k}\nonumber\\
&&\!\!\!\!\!\!
={\frac{1}{4{ \pi }^{2}}}
\int {\frac{
\exp \left({i\ {\bmi  k}
\cdot {\bmi  r}_{i}}\right) 
\exp \left({i\ {\bmi  k}
\cdot ({\bmi  R}_{i} -{\bmi  R}_{j})}\right)
\exp \left({-i\ {\bmi  k}
\cdot {\bmi  r}_{j}}\right)} 
{\hat{ K}_{c}^{2}-{k}^{2}}}\ {{\rm d}}{\bmi  k}
\end{eqnarray}
where  ${k}^{2}={k}_{z}^{2}+{k}_{x}^{2}$.
We now introduce a set of {\em shifted cylindrical coordinate\/} systems
so that
\begin{eqnarray} 
{\bmi  k}&=&k(\cos \lambda \ ,\sin \lambda)\nonumber\\ 
{ {\rm d}}{\bmi  k}&=&k {{\rm d}}{ k} {{\rm d}}{\lambda}\nonumber\\ 
{\bmi  r}_{i}&=&{\bmi  x}-{\bmi  R}_i=r_{i}(\cos \theta_i \ ,\sin 
\theta_i)\nonumber\\
{\bmi  R}_{ij}&=&{\bmi  R}_{i}-{\bmi  R}_{j}= R_{ij}(\cos \theta_{ij} \ 
,\sin\theta_{ij})\nonumber\\
 {\bmi  r}_{j}&=&{\bmi  x}^{\prime}-{\bmi  R}_{j}=r_{j}(\cos \theta_j \ ,\sin 
 \theta_j)
\end{eqnarray} 
hence (\ref{Ig1}) is transformed into
\begin{eqnarray} \label{Ig2} {g}_{c }({\bmi  x} -{\bmi  x}^{\prime} 
)&=&{\frac{1}{4{ \pi }^{2}}}
\int_{0}^{\infty}  
\frac{{{\rm d}}{ k}\  k}{{\hat{ K}_{c }^{2}-{k}^{2}}} \times \nonumber\\ &&
\int_{0}^{2 \pi}{{\rm d}}{\lambda}
\exp i k \left[ {  r_{i}}\cos(\theta_i-\lambda) + {  
R}_{ij}\cos(\theta_{ij}-\lambda) -  {
r}_{j} \cos(\theta_j-\lambda) 
\right] .
\end{eqnarray} 

Using the {\em Jacobi-Anger\/}  expansion [e.g {\em 
Arfken}~\shortcite{Arfken}, {\em p.}$585$]
\begin{equation} \label{JacobiAnger}
\exp(i r \cos s)=\sum\limits_{ n\ =\ -\infty }^{ \infty } { i}^{ n} \
{J}_{n}(r)\ \exp \left({i\ n\ s}\right)
\end{equation} for arbitrary real numbers $r$ and $s$, we transform the
exponential appearing in (\ref{Ig2}) into
\begin{eqnarray} \label{Ig3} 
\lefteqn{{g}_{c }({\bmi  x} -{\bmi  x}^{\prime} )=\frac{1}{{4 \pi }^{2}}
\sum\limits_{ m,n,p = -\infty }^{ \infty } (-1)^{n}{ i}^{ n+m+p}
\int_{0}^{\infty}  
\frac{{{\rm d}}{ k}\  k}{{\hat{ K}_{c }^{2}-{k}^{2}}} {J}_{m}(k r_i) \ {J}_{n}(k
r_{j})\ {J}_{p}(k  R_{ij})
}\nonumber\\  
&&\times \int_{0}^{2 \pi}{{\rm d}}{\lambda} \exp i  \left[
m(\theta_i-\lambda) + p (\theta_{ij}-\lambda ) -n(\theta_j-\lambda)
\right]\nonumber\\
&&={\frac{1}{2{ \pi }}}
\sum\limits_{ m,n,p = -\infty }^{ \infty } (-1)^{n}{ i}^{ n+m+p}
\delta_{n-m-p}   \exp i  \left[ m \theta_i + p \theta_{j} -n\theta_j 
\right]
\nonumber\\  
&&\times \int_{0}^{\infty}  
\frac{{{\rm d}}{ k}\  k}{{\hat{ K}_{c }^{2}-{k}^{2}}} {J}_{m}(k r_i) \ {J}_{n}(k
r_j)\ {J}_{p}(k  R_{ij})\nonumber\\
&&={\frac{1}{2{ \pi }}}
\sum\limits_{ m,n = -\infty }^{ \infty }  {\rm e}^{ i  \left[ m \theta_i + 
(n-m)
\theta_{ij} -n\theta_j \right]}\int_{0}^{\infty}  
\frac{{{\rm d}}{ k}\  k}{{\hat{ K}_{c }^{2}-{k}^{2}}} {J}_{m}(k r_i)  {J}_{n}(k
r_j) {J}_{n-m}(k  R_{ij}).
\end{eqnarray}


Assuming the geometrical (or {\em muffin tin\/}) constraint [see {\em 
Gonis\/}~(1992)]  
\begin{equation} \label{Muffin}
 R_{ij} > r_{i} + r_{j},
\end{equation}  
so that ${\bmi r}_{i}$ and ${\bmi r}_{j}$ are inside {\em non-overlapping\/} discs 
centred at ${\bmi R}_{i}$ and ${\bmi R}_{j}$ respectively,
we can evaluate the $k$-integral in (\ref{Ig3}) using the calculus of residues.
 The geometrical constraint 
(\ref{Muffin})
 is necessary to ensure the {\em uniform convergence\/} 
of subsequent contour integrals.

Given the {\em unit \/} cartesian $2-$vectors $\hat{\bmi k}=({k}_{z},{k}_{x})/k$ and 
$\hat{\bmi r}=({z},{x})/r$
where $k^{2}={k}_{z}^{2}+{k}_{x}^{2}$ and $r^{2}={z}^{2}+{x}^{2}$, 
we define the {\em cylindrical harmonics\/}
\begin{equation} \label{Ydef}
{Y}_{n}(\hat{\bmi k})={\left({{\frac{{k}_{z}+i{k}_{x}}{k}}}\right)}^{n}, 
\quad \mbox{and} \quad 
{Y}_{n}(\hat{\bmi r})={\left({{\frac{{z}+i{x}}{r}}}\right)}^{n}
\end{equation}
and their complex conjugates
\begin{equation} \label{Ycdef}
{Y}_{n}^{\ast}(\hat{\bmi k})={\left({{\frac{{k}_{z}-i{k}_{x}}{k}}}\right)}^{n}, 
\quad \mbox{and} \quad 
{Y}_{n}^{\ast}(\hat{\bmi r})={\left({{\frac{{z}-i{x}}{r}}}\right)}^{n}.
\end{equation}

Recalling the standard Bessel function  identity
${J}_{n-m}(k  R_{ij})=
\frac{1}{2}({H}^{(1)}_{n-m}(k  R_{ij})+ {H}^{(2)}_{n-m}(k  R_{ij}))$, and 
evaluating the $k$ integral in (\ref{Ig3}) using the residue theorem, we get
\begin{eqnarray} \label{Ig4} 
\lefteqn{
g_{c}({\bmi  x} -{\bmi  x}^{\prime} )}\nonumber\\
&&=\frac{-i}{4}
\sum\limits_{ m,n=-\infty }^{ \infty }  \exp i  \left[ m \theta_i + (n-m)
\theta_{ij} -n\theta_j \right]
 {J}_{m}(\hat{K}_{c} r_i) {J}_{n}(\hat{K}_{c} r_j) \
{H}^{(1)}_{n-m}(\hat{K}_{c}R_{ij})\nonumber\\
&&=
{\frac{-i}{4}} \sum\limits_{ n,m=- \infty }^{ \infty } 
{ J}_{ m} (\hat{ K}_{c}{r}_{i}){Y}_{m}({\hat{\bmi r}}_{i}){H}_{n-m}^{(1)}(\hat{ K}_{c}{R}_{ij})
{Y}_{n-m}(\hat{{\bmi R}}_{ij}){J}_{n}(\hat{ K}_{c}{r}_{j}){Y}_{n}^{{}^\ast}
({\hat{\bmi r}}_{j})\nonumber\\
&&=
{\frac{-i}{4}} \sum\limits_{ n,m=- \infty }^{ \infty } 
{ J}_{ m} (\hat{ K}_{c}{r}_{i}){Y}_{m}({\hat{\bmi r}}_{i}){H}_{m-n}^{(1)}(\hat{ K}_{c}{R}_{ij})
{Y}_{m-n}^{{}^\ast}(\hat{{\bmi R}}_{ji}){J}_{n}(\hat{ K}_{c}{r}_{j}){Y}_{n}^{{}^\ast}
({\hat{\bmi r}}_{j})
\end{eqnarray}

Upon using $g_c$ of (\ref{Ig4}) in the expression (\ref{Gdef1}) we find
the Green's tensor in terms of its {\em cylindrical partial wave\/} components, {\em viz.\/}
\begin{displaymath}
\fbox{
\begin{minipage} {143.0mm}
\begin{eqnarray} \label{G1Trep}
\lefteqn{\rf{\bmi G}({\bmi x},{\bmi x}^{\prime};k_y)=}\nonumber\\
&&\!\!\!\!\!\!\!\!\!\!\!\!
\frac{-i}{4 \rho \omega^2}\!\!\!\!
\sum\limits_{\stackrel{m,n}{c\in\{P,S,H\}}}\!\! \!\!\!\! 
{H}_{m-n}^{(1)}(\hat{ K}_{c}{R}_{ij})
{Y}_{m-n}^{{}^\ast}(\hat{{\bmi R}}_{ji})
{\bmi L}_{c}\left[{ J}_{ m} (\hat{ K}_{c}{r}_{i}){Y}_{m}({\hat{\bmi r}}_{i})  \right] \otimes 
{\bmi  L}_{c}^{\prime}\left[{J}_{n}(\hat{ K}_{c}{r}_{j}){Y}_{n}^{{}^\ast}
({\hat{\bmi r}}_{j})  \right]{\bmi J}_{6}\nonumber\\
 &&\!\!\!\!\!\!\!\!\!\!\!\!=\frac{-i}{4 \rho \omega^2}\!\!\!\!
\sum\limits_{\stackrel{m,n}{c\in\{P,S,H\}}} \!\!\!\!\!\!
{H}_{m-n}^{(1)}(\hat{ K}_{c}{R}_{ij})
{Y}_{m-n}^{{}^\ast}(\hat{{\bmi R}}_{ji})
 |(\re,k_y,c,m),{\bmi r}_{i}\rangle 
 \langle(\re,k_y,c,n),{\bmi r}_{j}|{\bmi J}_{6}
\end{eqnarray} 
\vspace{2mm}
\end{minipage}
}
\end{displaymath}
where
\[
{\bmi x}={\bmi R}_i+{\bmi x}_i, \quad {\bmi x}^{\prime}={\bmi R}_j+{\bmi 
x}_j,
\quad  R_{ij}=\|{\bmi R}_j-{\bmi R}_i\|,
\quad \hat{{\bmi R}}_{ji}= ({\bmi R}_j-{\bmi R}_i)/R_{ij},
\]
the $6-$component basis states 
$|(\re,k_y,c,m),{\bmi r}_{i}\rangle$ are defined in {\sc Box} {\bf 4.1}, 
with the dual states $\langle(\re,k_y,c,m),{\bmi r}_{i}|$ given by the rule
\[
\langle(\re,k_y,c,m),{\bmi r}_{i}|\equiv (-1)^m |(\re,-k_y,c,-m),{\bmi 
r}_{i}\rangle^{T}.
\]
%%%%%%%%%%%%%%%%%%%%%%%%%%%%%%%%%%%%%%%%%%%%%%%%%%%%%%%%%%%%%%%%%%%%%%%%%%%%%%%%%%%%%%%%%%%%
\begin{figure}
\fbox{
\begin{minipage} {140.0mm}
{\sc \bf  Box} ${\bf  4.1}$ {\sc The $6-$component cylindrical polar basis states}  \newline
Given that
\[
M=m(m+1), \qquad {Y}_{m}(\hat{\bmi r})=\frac{\left( z\, + i\,x \right)^m}{r^m}
\]

The {\em regular\/} basis states are defined as
\[
|(\re,k_y,c,m),{\bmi r}\rangle \equiv {\bmi L}_{c}
\left[ {J}_{m}(\hat{K}_{c}  r) {Y}_{m}(\hat{\bmi r}) \right], \quad c\in 
\{ P,S,H \}
\]
and the {\em outgoing\/} basis states are defined as
\[
|(\ou,k_y,c,m),{\bmi r}\rangle \equiv {\bmi L}_{c}
\left[ {H}^{(1)}_{m}(\hat{K}_{c}  r) {Y}_{m}(\hat{\bmi r}) \right], \quad c\in 
\{ P,S,H \}
\]
where the differential operators ${\bmi L}_{c}$ were defined in (\ref{Labdef}).

Then, assuming $C_m\in \{ {J}_{m}, {H}^{(1)}_{m} \}$ we have
\begin{samepage}
\begin{eqnarray} \label{BPdef}
\lefteqn{|(\{\re, \ou\},k_y,P,m),{\bmi r}\rangle \equiv {\bmi L}_{P}\left[ 
{C}_{m}(\hat{K}_{P}  r)
{Y}_{m}(\hat{\bmi r}) \right]=
{{Y}_{m}(\hat{\bmi r})}
\times}\nonumber\\
&&\!\!\!\!\!\!
\left[\begin{array}{c}
\frac{ C_{m-1}(\hat{K}_{P} r)\,\hat{K}_{P}\, z}{ r} - 
    \frac{C_{m}(\hat{K}_{P} r)\,m\,\left( z + i\,x \right) }{r^2}\\
{{C_{m-1}(\hat{K}_{P} r)\,\hat{K}_{P}\,x}\over {r}} -  
    {{{C_{m}(\hat{K}_{P} r)}\,m\,\left(   x\,- i\,z \right) }\over 
{{r^2}}}\\ 
i k_{y} {C_{m}(\hat{K}_{P} r)}\\
{{2\,\rf{\mu}\,C_{m-1}(\hat{K}_{P} r) \,\hat{K}_{P}\, 
        \left( {x^2} -{z^2} - 2\,i\,m\,z\,x  \right) }\over {{r^3}}} -  
{\scriptstyle{C_{m}(\hat{K}_{P} r)}}\left[ \lambda \, {K}^{2}_{P} -  
{{2\,\rf{\mu} \, 
           \left( -r^2\,z^2\,\hat{K}^{2}_{P}
+ M \,{{\left( z + i\,x \right) }^2} 
             \right) }\over 
          {{r^4}}} \right]\\
2 \rf{\mu} \,\left[{{C_{m-1}(\hat{K}_{P} r) \,\hat{K}_{P}\, 
        \left(i\,m\,({z^2} -\,{x^2})   - 2\,z\,x  \right) 
}\over{{r^3}}} 
- {{{C_{m}(\hat{K}_{P} r)} \, \left( 
r^2\,z\,x\,\hat{K}^{2}_{P} 
+  i\,M  \,{{\left( z + i\,x \right) }^2} \right) } \over {{r^4}}}\right]\\
 2\, \rf{\mu}\,i\, k_{y}\left[ \frac{ C_{m-1}(\hat{K}_{P} r)\,\hat{K}_{P}\, z}{ r} - 
    \frac{C_{m}(\hat{K}_{P} r)\,m\,\left( z\, + i\,x \right) }{r^2}\right]
\end{array}
\right];
\end{eqnarray}
\end{samepage}
\begin{samepage}
\begin{eqnarray} \label{BSdef}
\lefteqn{|(\{\re, \ou\},k_y,S,m),{\bmi r}\rangle \equiv {\bmi L}_{S}\left[ {C}_{m}(\hat{K}_{S}  r)
{Y}_{m}(\hat{\bmi r}) \right]=
\frac{\omega \,{{Y}_{m}(\hat{\bmi r})}}{\rf{\beta} \hat{K}_{S}}
\times}\nonumber\\
&&\!\!\!\!\!\!
\left[
\begin{array}{c}
{{C_{m-1}(\hat{K}_{S} r)\,\hat{K}_{S}\,x}\over {r}} - 
    {{{C_{m}(\hat{K}_{S} r)}\,m\,\left(   x\, - i\,z \right) }\over 
{{r^2}}}\\
-{{C_{m-1}(\hat{K}_{S} r)\,\hat{K}_{S}\,z}\over {r}} + 
    {{{C_{m}(\hat{K}_{S} r)}\,m\,\left( z\, + i\,x \right) }\over {{r^2}}}\\
    0\\
2\, \rf{\mu}\,\left[{{C_{m-1}(\hat{K}_{S} r) \,\hat{K}_{S}\,
        \left(i\,m\,\left( {z^2} - {x^2} \right) -2\,z\,x   \right) 
}\over 
      {{r^3}}}-
{{{C_{m}(\hat{K}_{S} r)} \,\left( 
r^2\,\,z\,x\,\hat{K}^{2}_{S} + 
          i\,M \,{{\left( z\, + i\,x \right) }^2} \right) }\
       \over {{r^4}}}\right]\\
2 \rf{\mu} \left[{{C_{m-1}(\hat{K}_{S} r) \,\hat{K}_{S}\,
        \left( {z^2} + 2\,i\,m\,z\,x - {x^2} \right) }\over {{r^3}}} + 
{C_{m}(\hat{K}_{S} r)}\,\left( \frac{\hat{K}^{2}_{S} (z^2 - x^2)}{2 r^2}
- {{ M \,{{\left( z\, + i\,x \right) }^2}  }
             \over {{r^4}}}\right)\right]\\ 
 \rf{\mu}\, i\, k_{y}\left[{{C_{m-1}(\hat{K}_{S} r)\,\hat{K}_{S}\,x}\over {r}} - 
    {{{C_{m}(\hat{K}_{S} r)}\,m\,\left(  x\, - i\,z \right) }\over 
{{r^2}}}\right]
\end{array}
\right];
\end{eqnarray}
\end{samepage}
\begin{samepage}
\begin{eqnarray} \label{BHdef}
\lefteqn{|(\{\re, \ou\},k_y,H,m),{\bmi r}\rangle \equiv {\bmi L}_{H}\left[ {C}_{m}(\hat{K}_{H}  r)
{Y}_{m}(\hat{\bmi r}) \right]=
\frac{{{Y}_{m}(\hat{\bmi r})}}{\hat{K}_{H}}
\times}\nonumber\\
&&\!\!\!\!\!\!
\left[\begin{array}{c} 
i\, k_{y}\left[\frac{ C_{m-1}(\hat{K}_{H} 
r)\,\hat{K}_{H}\, z}{ r} - \frac{C_{m}(\hat{K}_{H} r)\,m\,\left( 
z\, + i\,x \right) }{r^2}\right]\\ 
i\, k_{y} \left[{{C_{m-1}(\hat{K}_{H} 
r)\,\hat{K}_{H}\,x}\over {r}} - {{{C_{m}(\hat{K}_{H} r)}\,m\,\left( 
x\, - i\,z \right) }\over {{r^2}}}\right]\\ \hat{K}^{2}_{H} 
{C_{m}(\hat{K}_{H} r)}\\ 2\, \rf{\mu}\, i\, k_{y}\, \left[{{C_{m-1}(\hat{K}_{H} 
r) \,\hat{K}_{H}\, \left( {x^2} -{z^2} - 2\,i\,m\,z\,x \right) }\over 
{{r^3}}} + {{{C_{m}(\hat{K}_{H} r)}\, \left( 
-r^2\,{z^2}\,\hat{K}^{2}_{H} + M \,{{\left( z\, + i\,x \right) }^2} 
\right) }\over {{r^4}}} \right]\\ 
2\,\rf{\mu}\,i\, k_{y} \left[ 
{{C_{m-1}(\hat{K}_{H} r) \,\hat{K}_{H}\, \left(i\,m\,({z^2} 
-\,{x^2}) - 2\,z\,x \right) }\over{{r^3}}} - {{{C_{m}(\hat{K}_{H} r)} 
\, \left( r^2\,z\,x\,\hat{K}^{2}_{H} + i\,M \,{{\left( z\, + i\,x \right) 
}^2} \right) } \over {{r^4}}}\right]\\ \rf{\mu} \left(\hat{K}_{H}^{2}- 
k_{y}^{2} \right)\left[\frac{ C_{m-1}(\hat{K}_{H} r)\, 
\hat{K}_{H}\, z}{ r} - \frac{C_{m}(\hat{K}_{H} r)\,m\,\left( z + 
i\,x \right) }{r^2}\right] 
\end{array}
\right];
\end{eqnarray}
\end{samepage}
with dual basis vectors 
\[
\langle(\{\re, \ou\},k_y,c,m),{\bmi r}|=(-1)^m |(\{\re, \ou\},-k_y,c,-m),{\bmi r}\rangle^{T},\quad
c\in\{P,S,H\}  
\]
\vspace{6mm}
\end{minipage}
}
\end{figure}

%%%%%%%%%%%%%%%%%%%%%%%%%%%%%%%%%%%%%%%%%%%%%%%%%%%%%%%%%%%%%%%%%%%%%%%%%%%%%%%%%%%%%
%%%%%%%%%%%%%%%%%%%%%%%%%%%%%%%%%%%%%%%%%%%%%%%%%%%%%%%%%%%%%%%%%%%%%%%%%%%%%%%%%%%%%
%%%%%%%%%%%%%%%%%%%%%%%%%%%%%%%%%%%%%%%%%%%%%%%%%%%%%%%%%%%%%%%%%%%%%%%%%%%%%%%%%%%%%
%%%%%%%%%%%%%%%%%%%%%%%%%%%%%%%%%%%%%%%%%%%%%%%%%%%%%%%%%%%%%%%%%%%%%%%%%%%%%%%%%%%%%
%%%%%%%%%%%%%%%%%%%%%%%%%%%%%%%%%%%%%%%%%%%%%%%%%%%%%%%%%%%%%%%%%%%%%%%%%%%%%%%%%%%%%
%%%%%%%%%%%%%%%%%%%%%%%%%%%%%%%%%%%%%%%%%%%%%%%%%%%%%%%%%%%%%%%%%%%%%%%%%%%%%%%%%%%%%
%%%%%%%%%%%%%%%%%%%%%%%%%%%%%%%%%%%%%%%%%%%%%%%%%%%%%%%%%%%%%%%%%%%%%%%%%%%%%%%%%%%%%
%%%%%%%%%%%%%%%%%%%%%%%%%%%%%%%%%%%%%%%%%%%%%%%%%%%%%%%%%%%%%%%%%%%%%%%%%%%%%%%%%%%%%%%%%%%
%%%%%%%%%%%%%%%%%%%%%%%%%%%%%%%%%%%%%%%%%%%%%%%%%%%%%%%%%%%%%%%%%%%%%%%%%%%%%%%%%%%%%%%%%%%
%%%%%%%%%%%%%%%%%%%%%%%%%%%%%%%%%%%%%%%%%%%%%%%%%%%%%%%%%%%%%%%%%%%%%%%%%%%%%%%%%%%%%%%%%%%
%%%%%%%%%%%%%%%%%%%%%%%%%%%%%%%%%%%%%%%%%%%%%%%%%%%%%%%%%%%%%%%%%%%%%%%%%%%%%%%%%%%%%%%%%%%






\section{An effective point scattering model} \label{secEffScat}
In this section we will take care of the nearfield singularity of the 
reference medium Green's tensor ${\rf{\bmi G}}$ in the limit as the 
source point and field point coincide. The plane wave representations of 
the Green's tensor from chapter \ref{Grep} are ill suited to this task, 
and we will need to employ 
the {\em spatial domain\/} representation (\ref{G1Trep}) of ${\rf{\bmi G}}$. Using this 
newly acquired Green's tensor, we will sum the multiple scattering series 
for the special case of a uniform medium with heterogeneity supported 
exclusively on the square subregion
${\mit \Omega}_j$ with a midpoint at ${\bmi R}_j$, and 
a side length ${\mit \Delta}={\mit \Delta}x={\mit \Delta}z$ 
that is small on the scale of the incident 
wavelength. We let the square ${\mit \Omega}_j$ shrink to point, while 
maintaining its square shape. All of the singular multiple scattering 
interactions may then be summarized by replacing the material contrasts 
responsible for the scattering by their {\em effective \/} values, i.e. 
{\em a single scattering interaction with  the effective material contrasts 
on ${\mit \Omega}_j$ is equivalent to the sum of all multiple scattering 
interactions from the actual material contrasts\/}. This approach has two 
advantages:
\begin{itemize}
\item we replace the multiple scattering series with a 
single scattering interaction, and consequently, we avoid the nearfield 
singularity of  ${\rf{\bmi G}}$

\item the effective material contrasts are generally smaller than the 
actual contrasts. As a result the multiple scattering series appropriate 
to the  effective medium scattering problem is 
easier  to sum than the corresponding series in the original medium. This 
is consistent with the numerical experience of {\em Schuster\/}~(1985), 
who employed a similar renormalization strategy to solve a multi-inclusion 
acoustic scattering problem using a boundary integral equation method.
\end{itemize}

%%%%%%%%%%%%%%%%%%%%%%%%%%%%%%%%%%%%%%%%%%%%%%%%%%%%%%%%%%%%%%%%%%%%%%%%%%%%%%%%%%%%%%%
%%%%%%%%%%%%%%%%%%%%%%%%%%%%%%%%%%%%%%%%%%%%%%%%%%%%%%%%%%%%%%%%%%%%%%%%%%%%%%%%%%%%%%%
%%%%%%%%%%%%%%%%%%%%%%%%%%%%%%%%%%%%%%%%%%%%%%%%%%%%%%%%%%%%%%%%%%%%%%%%%%%%%%%%%%%%%%%
%%%%%%%%%%%%%%%%%%%%%%%%%%%%%%%%%%%%%%%%%%%%%%%%%%%%%%%%%%%%%%%%%%%%%%%%%%%%%%%%%%%%%%%
\subsection{Single scattering}

The {\em single scattering interaction\/} at ${\mit 
\Omega}_j$ may be expressed as 
\begin{equation} \label{Sscat}
{\rf{\bmi G}}{\mit \Delta} {\bmi A} {\rf{\bmi G}}\equiv
\int_{{\mit \Omega}_j} {\rf{\bmi G}}({\bmi x}_{R};{\bmi x};k_y)
{\mit \Delta} {\bmi A}({\bmi
x},\partial_{x};k_y) {\rf{\bmi G}}({\bmi x},{\bmi x}_{S};k_y) {\rm 
d}{\bmi x}.
\end{equation}
Using (\ref{G1Trep}) for the components of ${\rf{\bmi G}}$ in (\ref{Sscat}) 
gives
\begin{eqnarray} \label{Sscat1}
\lefteqn{{\rf{\bmi G}}{\mit \Delta} {\bmi A} {\rf{\bmi G}}=
\sum\limits_{\stackrel{m,n}{c\in\{P,S,H\}}} 
\sum\limits_{\stackrel{m^{\prime},n^{\prime}}{c^{\prime}\in\{P^{\prime},S^{\prime},H^{\prime}\}}}
\frac{- 1}{16 \rf{\rho}\!{}^2 \omega^4 }
 } \nonumber\\
&&\times
 |(\re,k_y,c,m),{\bmi r}_{R}\rangle {H}_{m-n}^{(1)}(\hat{ K}_{c}{R}_{Rj})
{Y}_{m-n}^{{}^\ast}(\hat{{\bmi R}}_{jR})\nonumber\\
 &&
\times \int_{{\mit \Omega}_j}  \langle(\re,k_y,c,n),{\bmi x}-{\bmi R}_{j}|{\bmi J}_{6}
{\mit \Delta} {\bmi A}({\bmi
x},\partial_{x};k_y)  |(\re,k_y,c^{\prime},m^{\prime}),{\bmi x}-{\bmi R}_{j}\rangle 
{\rm d}{\bmi x}\nonumber\\
&&
\times {H}_{m^{\prime}-n^{\prime}}^{(1)}(\hat{ K}_{c^{\prime}}{R}_{jS})
{Y}_{m^{\prime}-n^{\prime}}^{{}^\ast}(\hat{{\bmi R}}_{Sj}) 
 \langle(\re,k_y,c^{\prime},n^{\prime}),{\bmi r}_{S}|{\bmi J}_{6}
\end{eqnarray} 
where 
\[
{\bmi x}_{R}={\bmi R}_{R}+{\bmi r}_{R}, \qquad {\bmi x}_{S}={\bmi 
R}_{S}+{\bmi r}_{S}
\]
in (\ref{Sscat}) and (\ref{Sscat1}).
The integrals 
\begin{eqnarray} \label{Celems}
\lefteqn{{\mit \Delta}{\bmi C}^{j}_{c,n;c^{\prime},m^{\prime}}=\int_{{\mit \Omega}_j}
\left\langle (\re,k_y,c,n),{\bmi x}-{\bmi R}_{j}\left|
{\bmi J}_{6}{{\mit \Delta} {\bmi A}}({\bmi x};{\partial 
}_{x},{k}_{y})\right|(\re,k_y,c^{\prime},m^{\prime}),
{\bmi x}-{\bmi R}_{j}\right\rangle\, {\rm d}{\bmi x}}\nonumber\\
&&
\end{eqnarray}
in (\ref{Sscat1}) represent the {\em single scattering approximation\/} to the 
coupling between the multipoles induced by the material heterogeneity at 
${\bmi R}_{j}$. 


\begin{figure}
%\setlength{\fboxrule}{1mm}
\fbox{
\begin{minipage} {145.0mm}
{\sc  Box} ${\bf  4.2}$ {\sc Weak formulations for the displacement 
traction wavefield}\newline
The $L_{x}$-periodic wave field 
\[
{\bmi 
b}(z,x)={\left({\begin{array}{cccccc}{u}_{z}&{u}_{x}&{u}_{y}&
{ \tau }_{zz}&{ \tau}_{xz}&{ \tau}_{yz}\end{array}}\right)}^{T},
\] 
is continuous in a 
source free region, because it must satisfy the wave equation, so that 
the components are differentiable, and consequently, they are continuous. 
More explicitly,
\[
{u}_{z},{u}_{x},{u}_{y}\in C^{1}[0,L_x]\quad \mbox{and} \quad
{ \tau }_{zz},{ \tau}_{xz},{ \tau}_{yz}\in C^{0}[0,L_x].
\] 
However, we have not chosen to include ${ \tau }_{xx}$ in the dependent 
variable ${\bmi b}(z,x)$. Using a momentum balance argument, we know that
${ \tau }_{xx}$ should be continuous across transverse discontinuities of 
the medium properties. Yet, we have no guarantee that this condition will 
be fulfilled, because we do not model ${ \tau }_{xx}$ explicitly, instead 
we construct it  via the elastic wave stress-strain relations that may be 
expressed in the adopted notation as
\begin{equation} \label{TxxDef}
{\tau }_{xx}={\rf{\zeta}} {\partial }_{ x}{ u}_{ x}
+i\,{k}_{y} {{\rf{\chi}}}{ u}_{ y}+{\rf{\gamma}} {\tau }_{zz},
\end{equation}
where the material parameters ${\rf{\gamma}}$, ${\rf{\chi}}$ and ${\rf{\zeta}}$ where 
introduced in Table \ref{ATable} and are re-written here as
\[
{\rf{\gamma}} =\left[{1-2{\frac{\rf{\beta }{}^{2}}{\rf{\alpha}{}^{2}}}}\right]
\quad
{\rf{\chi}}=2 \rf{\rho} \rf{\beta}{}^2\left[1- 
2{\rf{\beta}{}^{2}\over\rf{\alpha}{}^{2}}\right]
\quad \mbox{ and}\quad 
{\rf{\zeta}}  =4 \rf{\rho} \rf{\beta }{}^{2}
\left[{ 1-{\frac{\rf{ \beta }{}^{2}}{\rf{ \alpha }{}^{2}}}}\right]
\]  

Now, let us assume  that the material parameters  ${\rf{\gamma}}$, ${\rf{\chi}}$ and ${\rf{\zeta}}$  in 
(\ref{TxxDef}) 
have a {\em finite jump discontinuity\/} at $x=d\in(0,L_x)$, so that
\[
\left[{ {\rf{\gamma}} }\right](d) \equiv   \lim\limits_{ 
\varepsilon    \rightarrow   0+}\  {\rf{\gamma}}   (d+ 
\varepsilon   )- {\rf{\gamma}}   (d- \varepsilon)\quad 
 \left[{ {\rf{\chi}} }\right](d) \equiv   \lim\limits_{ 
\varepsilon    \rightarrow   0+}\  {\rf{\chi}}   (d+ 
\varepsilon   )- {\rf{\chi}}   (d- \varepsilon)\]
\[
\mbox{and}\quad
\left[{ {\rf{\zeta}} }\right](d) \equiv   \lim\limits_{ 
\varepsilon    \rightarrow   0+}\  {\rf{\zeta}}   (d+ 
\varepsilon   )- {\rf{\zeta}}   (d- \varepsilon)
\] 
then the traction component ${\tau }_{xx}$ suffers a similar saltus, since
\[
\lim\limits_{ 
\varepsilon    \rightarrow   0+}\  {\tau }_{xx}   (d+ 
\varepsilon   )- {\tau }_{xx}   (d- \varepsilon)=
\left[{ {\rf{\zeta}} }\right]\!(d)\,{\partial }_{ x}{ u}_{ x}
+i\, k_y\,\left[{ {\rf{\chi}} }\right]\!(d)\,{ u}_{ y}
+\left[{ {\rf{\gamma}} }\right]\!(d)\, {\tau }_{zz}
\]
where we have used the continuity of $\partial_x{ u}_{ x}$, ${ u}_{y}$  
and ${\tau }_{zz}$, because 
these are components of the solution to a source free wave equation.

Given the $L_x$-periodic test function $\phi(x)$, we want to find the {\em 
projection\/}
\begin{equation} \label{projection}
\int_{0}^{L_x} \phi \partial_{x}{\tau }_{xx} { d}x\equiv\int_{0}^{d-} \phi \partial_{x}{\tau }_{xx} 
{ d}x+\int_{d+}^{L_x} \phi \partial_{x} {\tau }_{xx} {\rm d}x
\end{equation} 
and the right hand side of (\ref{projection}) follows from the left because 
the integral $\int_{d-}^{d+} (\cdot) {\rm d}x$ has a 
vanishing measure.


The projection (\ref{projection}) is mathematically equivalent to
\begin{equation} \label{projection1}
\int_{0}^{L_x} \phi \partial_{x}{\tau }_{xx} {\rm d}x\equiv -\int_{0}^{d-} 
(\partial_{x}\phi) {\tau }_{xx} 
{\rm d}x-\int_{d+}^{L_x} (\partial_{x}\phi)  {\tau }_{xx} {\rm d}x
+\left( \phi {\tau }_{xx}\right]_{0}^{d-}+\left( \phi {\tau 
}_{xx}\right]_{d+}^{L_x}
\end{equation} 
which follows after an integration by parts. Due to periodicity, 
$(\phi  {\tau }_{xx})(0)=(\phi  {\tau 
}_{xx})(L_x)$, so that (\ref{projection1}) may be written in the equivalent 
form
\begin{equation} \label{projection2}
\int_{0}^{L_x} \phi \partial_{x}{\tau }_{xx} {\rm d}x\equiv -\int_{0}^{L_x} 
(\partial_{x}\phi) {\tau }_{xx} 
{\rm d}x -\left[ \phi {\tau }_{xx}\right](d)
\end{equation} 
where

\vspace{3mm}
\end{minipage}
}
\end{figure}
\begin{figure}
\fbox{
\begin{minipage} {145.0mm}
{\sc \bf  Box} ${\bmi  4}$.${\bmi 2}${\em continued\/} \newline
$
\left[ \phi {\tau }_{xx}\right](d)\equiv \lim\limits_{ 
\varepsilon    \rightarrow   0+}\  \left(\phi {\tau }_{xx}\right)(d+ \varepsilon   )- 
\left(\phi {\tau }_{xx}\right)(d- \varepsilon   )
$\newline
If we now drop the jump term $\left[ \phi {\tau }_{xx}\right](d)$ from 
(\ref{projection2}) so that
\begin{equation} \label{projection3}
\int_{0}^{L_x} \phi \partial_{x}{\tau }_{xx} {\rm d}x\equiv -\int_{0}^{L_x} 
(\partial_{x}\phi) {\tau }_{xx} 
{\rm d}x
\end{equation} 
we are implicitly assuming that
\begin{equation} \label{weakjump}
\left[ \phi {\tau }_{xx}\right](d)\equiv 
\int_{ 0}^{ L_x}\phi(x) \delta 
(x-d)\left[{{\tau}_{xx}}\right](d)\,{\rm d}x\equiv 0
\end{equation}

for all periodic test functions $\phi(x)$. The identity 
(\ref{projection3}) is then a statement that ${\tau}_{xx}$ is continuous 
across jumps in $\rf{\gamma}$, $\rf{\chi}$ and $\rf{\zeta}$ in a {\em projectional sense\/}, i.e. 
${\tau}_{xx}$ is {\em weakly continuous\/} on the interval $x\in [-,L_x)$.
To achieve a better understanding of the {\em weak\/} nature of 
(\ref{weakjump}), we should contrast it with the analogous {\em strong form\/}   
\[
\| \left[{\tau }_{xx}\right](d) \|\equiv  0
\]
in particular, if we choose the {\em uniform norm\/} 
$\| \cdot  \|_{\infty}$, then
\begin{equation} \label{strongjump1}
\| \left[ {\tau }_{xx}\right](d) \|_{\infty}\equiv  \left[ {\tau }_{xx}\right](d)=0
\end{equation}
so that continuity is enforced {\em pointwise\/} at $x=d$. While 
(\ref{strongjump1}) $\Rightarrow$(\ref{weakjump}), the converse is not 
generally true, i.e. (\ref{strongjump1})$\, \not \! \Leftarrow$ (\ref{weakjump}).
In the method of this thesis, the wave equations are being solved in a 
weak sense, because the boundary conditions (\ref{Geqntotal}) on the  Green's  tensor, 
were enforced by projection onto the space of bandlimited, 
periodic wavefields in the reference medium. Consequently, it is logically consistent to enforce 
the continuity of ${\tau }_{xx}$ in a projectional sense also.


Finally, the weak approximation of the {\em coupling integral\/} 
\[
\int {\rm d}x\left[{\begin{array}{cccccc}
{\tilde{u}}_{z}&{\tilde{u}}_{x}&{\tilde{u}}_{y}&{\tilde{ 
\tau }}_{zz}&{\tilde{ \tau }}_{xz}&{\tilde{ \tau 
}}_{yz}\end{array}}\right]{\bmi J}_6 \times\]
\[
\left[{\begin{array}{cccccc}
0&- {{\mit \Delta} \gamma} {\partial }_{  x}&  
-i {{\mit \Delta} \gamma} {  k}_{  y}&  {{\mit \Delta} a}&0&0\\0&0&0&0&{{\mit \Delta} b}&0\\0&0&0&0&0&{{\mit \Delta} b}\\ - {{\mit \Delta} \rho} {\omega }^{  
2}&  0&0&0&0&0\\ 0&- {{\mit \Delta} \rho} {\omega 
}^{  2}  -{ \partial }_{x} {{\mit \Delta} \zeta} {\partial }_{  
x}  + {{\mit \Delta} \mu} {  k}_{  y}^{  2}&  
-i{k}_{y}{ \partial }_{x} {{\mit \Delta} \chi}   -i{k}_{y} {{\mit \Delta} \mu} {\partial 
}_{  x}&  -{ \partial }_{x} {{\mit \Delta} \gamma} &  0&0\\ 
0&-i{k}_{y} {{\mit \Delta} \chi} {\partial }_{  x}  -i{k}_{y}{ \partial 
}_{x} {{\mit \Delta} \mu} &  - {{\mit \Delta} \rho} {\omega }^{  2}  + {{\mit \Delta} \zeta }
{  k}_{  y}^{  2}  -{ \partial }_{x} {{\mit \Delta} \mu} 
{\partial }_{  x}&  -i{k}_{y} {{\mit \Delta} \gamma} &  0&0\end{array}}\right]
\left[{\begin{array}{c}{u}_{z}\\ 
{u}_{x}\\
{u}_{y}\\
{ \tau }_{zz}\\
{ \tau }_{xz}\\
{ \tau }_{yz}\end{array}}\right]
\]
may be written as
\[
\int {\rm d}x\left[{\begin{array}{cccccccc}
{\tilde{u}}_{z}&{\tilde{u}}_{x}&{\tilde{u}}_{y}&{\tilde{ 
\tau }}_{zz}&{\tilde{ \tau }}_{xz}&{\tilde{ \tau 
}}_{yz}&\partial_x {\tilde{u}}_{x}&\partial_x {\tilde{u}}_{y}\end{array}}\right]\times\]
\[\left[{\begin{array}{cccccccc}
\displaystyle - {\mit \Delta}\rho {\omega }^{ 2}& \displaystyle 0& 
\displaystyle 0& \displaystyle 0& \displaystyle 0& \displaystyle 0& 
\displaystyle 0& \displaystyle 0\\ \displaystyle 0& \displaystyle {\mit \Delta}\mu { 
k}_{ y}^{ 2} - {\mit \Delta}\rho {\omega }^{ 2}& \displaystyle 0& \displaystyle 0& 
\displaystyle 0& \displaystyle 0& \displaystyle 0& \displaystyle -i{k}_{y} 
{\mit \Delta}\mu \\ \displaystyle 0& \displaystyle 0& \displaystyle {\mit \Delta}\zeta { k}_{ y}^{ 2} 
- {\mit \Delta}\rho {\omega }^{ 2}& \displaystyle -i{k}_{y} {\mit \Delta}\gamma & \displaystyle 0& 
\displaystyle 0& \displaystyle -i{k}_{y} \chi & \displaystyle 0\\ 
\displaystyle 0& \displaystyle 0& \displaystyle i {\mit \Delta}\gamma { k}_{ y}& 
\displaystyle -{\mit \Delta}a& \displaystyle 0& \displaystyle 0& \displaystyle {\mit \Delta}\gamma & 
\displaystyle 0\\ \displaystyle 0& \displaystyle 0& \displaystyle 0& 
\displaystyle 0& \displaystyle -{\mit \Delta}b& \displaystyle 0& \displaystyle 0& 
\displaystyle 0\\ \displaystyle 0& \displaystyle 0& \displaystyle 0& 
\displaystyle 0& \displaystyle 0& \displaystyle -{\mit \Delta}b& \displaystyle 0& 
\displaystyle 0\\ \displaystyle 0& \displaystyle 0& \displaystyle i{k}_{y} 
{\mit \Delta}\chi & \displaystyle {\mit \Delta}\gamma & \displaystyle 0& \displaystyle 0& 
\displaystyle {\mit \Delta}\zeta & \displaystyle 0\\ \displaystyle 0& \displaystyle 
i{k}_{y} {\mit \Delta}\mu & \displaystyle 0& \displaystyle 0& \displaystyle 0& 
\displaystyle 0& \displaystyle 0& \displaystyle {\mit \Delta}\mu 
\end{array}}\right]
\left[\begin{array}{c}
{u}_{z}\\ 
{u}_{x}\\
{u}_{y}\\
{ \tau }_{zz}\\
{ \tau }_{xz}\\
{ \tau }_{yz}\\
\partial_x{u}_{x}\\
\partial_x{u}_{y}\end{array}\right]
\]
\vspace{3mm}
\end{minipage}
}
\end{figure}
\begin{figure}
\fbox{
\begin{minipage} {145.0mm}
{\sc \bf  Box} ${\bmi  4}$.${\bmi 2}${\em continued\/} \newline
where
\[
\left[{\begin{array}{cccccc}
{{u}}_{z},&{{u}}_{x},&{{u}}_{y},&{{ 
\tau }}_{zz},&{{ \tau }}_{xz},&{{ \tau 
}}_{yz}\end{array}}\right]^{T}
\quad
\mbox{and }
\quad
\left[{\begin{array}{cccccc}
{\tilde{u}}_{z},&{\tilde{u}}_{x},&{\tilde{u}}_{y},&{\tilde{ 
\tau }}_{zz},&{\tilde{ \tau }}_{xz},&{\tilde{ \tau 
}}_{yz}\end{array}}\right]^{T}
\]
are the displacement/traction components of the incident and scattered 
partial wave states respectively.

\vspace{3mm}
\end{minipage}
}
\end{figure}
%%%%%%%%%%%%%%%%%%%%%%%%%%%%%%%%%%%%%%%%%%%%%%%%%%%%%%%%%%%%%%%%%%%%%%%%%%%%%%%%%%
%%%%%%%%%%%%%%%%%%%%%%%%%%%%%%%%%%%%%%%%%%%%%%%%%%%%%%%%%%%%%%%%%%%%%%%%%%%%%%%%%%


We evaluate the coupling integrals (\ref{Celems}) in the following two step process:
\begin{itemize}
\item we interpret the $\partial_x$ operators in ${{\mit \Delta} {\bmi A}}({\bmi x};{\partial 
}_{x},{k}_{y})$ in a {\em weak\/} or {\em projectional\/} fashion ({\em 
see\/} {\sc  Box} ${\bf  4.2}$). Consequently we employ integration by parts to transfer the 
$\partial_x$ operators from
${{\mit \Delta} {\bmi A}}({\bmi x};{\partial 
}_{x},{k}_{y})$ onto the basis vectors $\langle (\re,k_y,c,n),{\bmi 
x}-{\bmi R}_{j}|$ and  $|(\re,k_y,c^{\prime},m^{\prime}),
{\bmi x}-{\bmi R}_{j}\rangle$. The basis vectors are {\em augmented\/} by 
their partially differentiated components $(\partial_x u_x, \partial_x 
u_y)$. These augmented basis states are denoted by  
$\left|(\re,k_y,c,m),{\bmi r}\right)$,
where
\[
\left|(\re,k_y,c,m),{\bmi r}\right) \equiv \hat{\bmi L}_{c}\left[ 
{J}_{m}(\hat{K}_{c}  r)
{Y}_{m}(\hat{\bmi r}) \right], \ c\in\{P,S,H\}
\]
with
\begin{eqnarray} \label{Labaugdef}
\begin{array}{lll}
\hat{\bmi L}_{P}=&\hat{\bmi L}_{S}=&\hat{\bmi L}_{H}=\\
\left[\begin{array}{c}
\partial_{z} \\
\partial_{x} \\ 
\partial_{y} \\ 
2 \rf{\mu} \partial_{zz} - {\rf{\lambda} \omega^{2}\over \rf{\alpha}\!{}^2}\\
    2 \rf{\mu} \partial_{xz} \\ 
   2 \rf{\mu} \partial_{yz}\\
   \partial_{xx} \\
\partial_{xy}
\end{array}\right], \quad&
\frac{\omega }{\rf{\beta} \hat{K}_{S}}\,
\left[\begin{array}{c}
\partial_{x} \\-\partial_{z} \\ 0 \\
2 \rf{\mu} \partial_{xz}\\  
   \rf{\mu} (-\partial_{zz} + \partial_{xx}) \\ 
   \rf{\mu} \partial_{xy} \\
  - \partial_{xz} \\0
\end{array}\right], \quad&
\frac{1 }{\hat{K}_{H}}\,
\left[\begin{array}{c}
   \partial_{yz} \\
\partial_{xy} \\ 
{\omega^{2}\over \rf{\beta}\!{}^2} + \partial_{yy} \\ 
 2 \rf{\mu} \partial_{yzz}\\
    2 \rf{\mu} \partial_{xyz} \\ 
   \rf{\mu} ({\omega^{2}\over \rf{\beta}\!{}^2}\partial_{z} + 2 \partial_{yyz})\\
      \partial_{xxy} \\
{\omega^{2}\over \rf{\beta}\!{}^2}\partial_{x} + \partial_{xyy}  
\end{array}\right]
\end{array}
\end{eqnarray}
so that the $6-$vectors $\langle 
\cdot |$ and $| 
\cdot \rangle$ all concatenated with two extra components. 
The corresponding dual states $\left((\re,k_y,c,m),{\bmi r}\right|$ are 
given by the familiar rule
\[
\left((\re,k_y,c,m),{\bmi r}\right|\equiv (-1)^m \left|(\re,-k_y,c,-m),{\bmi 
r}\right)^{T}
\]
 
The resulting {\em weak\/} approximation 
to the elements of the coupling matrix (\ref{Celems}) may now be written 
as
\begin{samepage}
\begin{eqnarray} \label{weakRay0}
\lefteqn{
\int_{{\mit \Omega}_j}
\left\langle (\re,k_y,c,n),{\bmi x}-{\bmi R}_{j}\left|
{\bmi J}_{6}{{\mit \Delta} {\bmi A}}({\bmi x};{\partial 
}_{x},{k}_{y})\right|(\re,k_y,c^{\prime},m^{\prime}),
{\bmi x}-{\bmi R}_{j}\right\rangle\, {\rm d}{\bmi x}
}\nonumber\\
&&\qquad =\int_{{\mit \Omega}_j}
\left( (\re,k_y,c,n),{\bmi x}-{\bmi R}_{j}\left|
{{\mit \Delta} {\bmi C}}({\bmi x};{k}_{y})\right|(\re,k_y,c^{\prime},m^{\prime}),
{\bmi x}-{\bmi R}_{j}\right)\, {\rm d}{\bmi x}
\end{eqnarray}
\end{samepage}
where the {\em augmented coupling matrix\/} is defined as
\begin{samepage}
\begin{eqnarray}\label{CdefAug}
\lefteqn{
{\mit \Delta}{\bmi C}({\bmi R}_j;k_y)=}\\
&&\!\!\!\!\!\!
\left[{\begin{array}{cccccccc}
\displaystyle - {\mit \Delta}\rho {\omega }^{ 2}& \displaystyle 0& 
\displaystyle 0& \displaystyle 0& \displaystyle 0& \displaystyle 0& 
\displaystyle 0& \displaystyle 0\\ \displaystyle 0& \displaystyle {\mit \Delta}\mu { 
k}_{ y}^{ 2} - {\mit \Delta}\rho {\omega }^{ 2}& \displaystyle 0& \displaystyle 0& 
\displaystyle 0& \displaystyle 0& \displaystyle 0& \displaystyle -i{k}_{y} 
{\mit \Delta}\mu \\ \displaystyle 0& \displaystyle 0& \displaystyle {\mit \Delta}\zeta { k}_{ y}^{ 2} 
- {\mit \Delta}\rho {\omega }^{ 2}& \displaystyle -i{k}_{y} {\mit \Delta}\gamma & \displaystyle 0& 
\displaystyle 0& \displaystyle -i{k}_{y} \chi & \displaystyle 0\\ 
\displaystyle 0& \displaystyle 0& \displaystyle i {\mit \Delta}\gamma { k}_{ y}& 
\displaystyle -{\mit \Delta}a& \displaystyle 0& \displaystyle 0& \displaystyle {\mit \Delta}\gamma & 
\displaystyle 0\\ \displaystyle 0& \displaystyle 0& \displaystyle 0& 
\displaystyle 0& \displaystyle -{\mit \Delta}b& \displaystyle 0& \displaystyle 0& 
\displaystyle 0\\ \displaystyle 0& \displaystyle 0& \displaystyle 0& 
\displaystyle 0& \displaystyle 0& \displaystyle -{\mit \Delta}b& \displaystyle 0& 
\displaystyle 0\\ \displaystyle 0& \displaystyle 0& \displaystyle i{k}_{y} 
{\mit \Delta}\chi & \displaystyle {\mit \Delta}\gamma & \displaystyle 0& \displaystyle 0& 
\displaystyle {\mit \Delta}\zeta & \displaystyle 0\\ \displaystyle 0& \displaystyle 
i{k}_{y} {\mit \Delta}\mu & \displaystyle 0& \displaystyle 0& \displaystyle 0& 
\displaystyle 0& \displaystyle 0& \displaystyle {\mit \Delta}\mu 
\end{array}}\right]\nonumber 
\end{eqnarray}
\end{samepage}
with all material contrast parameters evaluated at the mid point ${\bmi 
R}_j$ of ${\mit \Omega}_j$. The operator is {\em multiplicative\/}, with the material contrast parameters in 
(\ref{CdefAug}) defined (with reference to Table \ref{ATable}) by Table \ref{contrastTable}, so that 
${\mit \Delta }\zeta$ is the difference between $\zeta$ inside the inclusion and 
the value of $\zeta$ in the homogeneous background medium,  with 
analogous definitions for the remaining contrasts. 



\item we evaluate the resulting integrals (\ref{weakRay0}) using the {\em midpoint\/} or 
{\em Rayleigh\/}  
approximation, to get the {\em bilinear form\/}
\begin{eqnarray} \label{weakRay}
\lefteqn{
\int_{{\mit \Omega}_j}
\left( (\re,k_y,c,n),{\bmi x}-{\bmi R}_{j}\left|
{\bmi J}_{6}{{\mit \Delta} {\bmi A}}({\bmi x};{\partial 
}_{x},{k}_{y})\right|(\re,k_y,c^{\prime},m^{\prime}),
{\bmi x}-{\bmi R}_{j}\right)\, {\rm d}{\bmi x}
}\nonumber\\
&&\qquad\qquad\qquad ={\mit \Delta}^{2}  \left( (\re,k_y,c,n),{\bf 0}+\left|
{\mit \Delta}{\bmi C}({\bmi R}_j;k_y)
\right|(\re,k_y,c^{\prime},m^{\prime}),{\bf 0}+\right)
\end{eqnarray}
where the {\em the augmented basis vectors\/} are given in {\sc Box} {\bf  
4.3} and we have used the Taylor series 
representations [e.g. {\em Ben-Menahem \& Singh\/}~(1981)]
\[
{J}_{n}(\hat{K}_{c}r)\underline{\sim 
}\left\{{\begin{array}{cc}
1-{\frac{{\left({\hat{K}_{c}r}\right)}^{2}}{2}}&n=0\\ 
{\frac{{\left({\hat{K}_{c}r}\right)}^{n}}{{2}^{n}n!}}&n>0
\end{array}}\right.
\quad \mbox{and employed}  \quad
{J}_{-n}(\hat{K}_{c}r)=(-1)^n {J}_{n}(\hat{K}_{c}r)
\]
where necessary.
\end{itemize}

%%%%%%%%%%%%%%%%%%%%%%%%%%%%%%%%%%%%%%%%%%%%%%%%%%%%%%%%%%%%%%%%%%%%%%%%%%%%%%%%%%%%%%%%
%%%%%%%%%%%%%%%%%%%%%%%%%%%%%%%%%%%%%%%%%%%%%%%%%%%%%%%%%%%%%%%%%%%%%%%%%%%%%%%%%%%%%%%%
%%%%%%%%%%%%%%%%%%%%%%%%%%%%%%%%%%%%%%%%%%%%%%%%%%%%%%%%%%%%%%%%%%%%%%%%%%%%%%%%%%%%%%%%
%%%%%%%%%%%%%%%%%%%%%%%%%%%%%%%%%%%%%%%%%%%%%%%%%%%%%%%%%%%%%%%%%%%%%%%%%%%%%%%%%%
%%%%%%%%%%%%%%%%%%%%%%%%%%%%%%%%%%%%%%%%%%%%%%%%%%%%%%%%%%%%%%%%%%%%%%%%%%%%%%%%%%
\begin{figure}
\fbox{
\begin{minipage} {145.0mm}
{\sc \bf  Box} ${\bmi  4}$.${\bmi 3}$ {\sc Multipole amplitudes of the 
augmented cylindrical harmonic basis states} \newline

When we take the limits
\[
\lim\limits_{{\bmi  x} \rightarrow {\bmi  R}_{j}}\ \left|(\re,k_y,c,m),{\bmi  x}- {\bmi  
R}_{j}\right)\equiv
\lim\limits_{r \rightarrow 0+}\ \hat{\bmi L}_{c}\left[ 
{J}_{m}(\hat{K}_{c}  r)
{Y}_{m}(\hat{\bmi r}) \right]
 \quad 
c\in\{P,S,H\}
\] 
 we find that:
\begin{itemize}
\item there are only $5$ nonvanishing  
harmonics $m\in \{-2,-1,0,1,2\}$ for the $P-$mode
\item there are only $4$ nonvanishing  
harmonics $m\in \{-2,-1,1,2\}$ for the $SV-$mode
\item there are only $5$ nonvanishing  
harmonics $m\in \{-2,-1,0,1,2\}$ for the $SH-$mode
\end{itemize}
We assemble these nonvanishing states $(\re,k_y,c,m), c\in \{P,S,H \}$ 
 into the $8 \times 14$ array 
$|\re,k_y,{\bf 0+})$  with the ordering
\begin{equation} \label{Bab0Def} 
|\re,k_y,{\bf 0+})=
\left[{\begin{array}{ccc}|(\re,k_y,P),{\bf 0+})&
|(\re,k_y,S),{\bf 0+})&
|(\re,k_y,H),{\bf 0+})\end{array}}\right]   
\end{equation}
where
\begin{equation} \label{P0rep}
\! |(\re,k_y,P),{\bf 0+})=\!\!\!\!
\begin{array}{ccccc}
\scriptstyle|(k_y,P,-2),{\bf 0}+)&\scriptstyle|(k_y,P,-1),{\bf 0}+)&
\scriptstyle|(k_y,P,0),{\bf 0}+)&
\scriptstyle|(k_y,P,1),{\bf 0}+)&\scriptstyle|(k_y,P,2),{\bf 0}+)\\
\left[\begin{array}{c} 0\\0\\0\\{{{\hat{K}_{P}^2}\,\rf{\mu}}\over 2}\\
  {-i\,\hat{K}_{P}^2\,\rf{\mu} \over 2}\\0\\{{-{\hat{K}_{P}^2}}\over 4}\\0
\end{array}\right.&
\begin{array}{c} {{-\hat{K}_{P}}\over 2}\\{i\,\hat{K}_{P}\over 2}\\0\\0\\0\\
-i\,k_y\,\hat{K}_{P}\,\rf{\mu}\\
  0\\{{- k_y\,\hat{K}_{P} }\over 2}
\end{array}&
\begin{array}{c} 0\\0\\i\,k_y\\- {{\omega^2}\,\lambda \over \rf{\alpha}{}^2}   - 
   {\hat{K}_{P}^2}\,\rf{\mu}\\0\\0\\{{-{\hat{K}_{P}^2}}\over 2}\\0
\end{array}&
\begin{array}{c} {{K}_{P}\over 2}\\{i\over 2}\,\hat{K}_{P}\\0\\0\\0\\i\,k_y\,\hat{K}_{P}\,\rf{\mu}\\0\\
  {{- k_y\,\hat{K}_{P} }\over 2}
\end{array}&
\left.
\begin{array}{c} 0\\0\\0\\{{{\hat{K}_{P}^2}\,\rf{\mu}}\over 2}\\
{i\,{\hat{K}_{P}^2}\,\rf{\mu}\over 2}\\0\\
  {{-{\hat{K}_{P}^2}}\over 4}\\0
\end{array}\right]\!\!
\end{array}
\end{equation}
\begin{equation} \label{S0rep}
|(\re,k_y,S),{\bf 0+})=\frac{\omega}{\rf{\beta}}\!\!
\begin{array}{cccc}
\scriptstyle|(k_y,S,-2),{\bf 0}+)&\scriptstyle|(k_y,S,-1),{\bf 0}+)&
\scriptstyle|(k_y,S,1),{\bf 0}+)&\scriptstyle|(k_y,S,2),{\bf 0}+)\\
\left[\begin{array}{c} 0\\0\\0\\{{-i \,{\hat{K}_{S}}\,\rf{\mu}}\over 2}\\
  {{-{\hat{K}_{S}}\,\rf{\mu} }\over 2}\\0\\{i\,{\hat{K}_{S}}\over 4}\\0
\end{array}\right.&
\begin{array}{c} {i\over 2}\\{1\over 2}\\0\\0\\0\\
  {{- k_y\,\rf{\mu}  }\over 2}\\0\\0
\end{array}&
\begin{array}{c} {i\over 2}\\{{-1}\over 2}\\0\\0\\0\\
  {{- k_y\,\rf{\mu}  }\over 2}\\0\\0
\end{array}&\left.
\begin{array}{c} 0\\0\\0\\{i\,{\hat{K}_{S}}\,\rf{\mu}\over 2}\\
  {{-{\hat{K}_{S}}\,\rf{\mu} }\over 2}\\0\\
  {{-i\,{\hat{K}_{S}}}\over 4}\\0
\end{array}\right]
\end{array}
\end{equation}
\begin{samepage}  
\begin{eqnarray} \label{H0rep}
\lefteqn{
|(\re,k_y,H),{\bf 0+})=}\\
&&\begin{array}{ccccc}
\scriptstyle|(k_y,H,-2),{\bf 0}+)&\scriptstyle|(k_y,H,-1),{\bf 0}+)&
\scriptstyle|(k_y,H,0),{\bf 0}+)&
\scriptstyle|(k_y,H,1),{\bf 0}+)&\scriptstyle|(k_y,H,2),{\bf 0}+)\\
\left[\begin{array}{c} 0\\0\\0\\{i\,k_y\,{\hat{K}_{H}}\,\rf{\mu}\over 2}\\
  {{k_y\,{\hat{K}_{H}}\,\rf{\mu}}\over 2}\\0\\
  {{-i\,k_y\,{\hat{K}_{H}}}\over 4}\\0
\end{array}\right.&
\begin{array}{c} {{-i\,k_y}\over 2}\\
{{- k_y }\over 2}\\0\\
  0\\0\\{{\rf{\mu}\,\left(2\, \rf{\beta}{}^2\, k_y^2-{\omega^2}  \right) }\over { 2\, \rf{\beta}{}^2} } 
  \\0\\{i\,{\hat{K}_{H}^2}\over 2}
\end{array}&
\begin{array}{c} 0\\0\\ \scriptstyle{\hat{K}_{H}}\\ \scriptstyle-i\,k_y\,{\hat{K}_{H}}\,\rf{\mu}\\0\\0\\
  {{-i\,k_y\,{\hat{K}_{H}}}\over 2}\\0
\end{array}&
\begin{array}{c} {i\,k_y\over 2}\\{{- k_y }\over 2}\\0\\0\\0\\
  {{\rf{\mu}\,\left({\omega^2}-2\, \rf{\beta}{}^2\, k_y^2  \right) }\over { 2\, \rf{\beta}{}^2} }\\
  0\\{i\,{\hat{K}_{H}^2}\over 2}
\end{array}&\left.
\begin{array}{c} 0\\0\\0\\{i\,k_y\,{\hat{K}_{H}}\,\rf{\mu}\over 2}\\
  {{- k_y\,{\hat{K}_{H}}\,\rf{\mu} }\over 2}\\0\\
  {{-i\,k_y\,{\hat{K}_{H}}}\over 4}\\0
\end{array}\right]
\end{array}\nonumber
\end{eqnarray}
\end{samepage}


%\vspace{3mm}
\end{minipage}
}
\end{figure}


%%%%%%%%%%%%%%%%%%%%%%%%%%%%%%%%%%%%%%%%%%%%%%%%%%%%%%%%%%%%%%%%%%%%%%%%%%%%%%%%%%%%%%
%%%%%%%%%%%%%%%%%%%%%%%%%%%%%%%%%%%%%%%%%%%%%%%%%%%%%%%%%%%%%%%%%%%%%%%%%%%%%%%%%%%%%%
%%%%%%%%%%%%%%%%%%%%%%%%%%%%%%%%%%%%%%%%%%%%%%%%%%%%%%%%%%%%%%%%%%%%%%%%%%%%%%%%%%%%%%
%%%%%%%%%%%%%%%%%%%%%%%%%%%%%%%%%%%%%%%%%%%%%%%%%%%%%%%%%%%%%%%%%%%%%%%%%%%%%%%%%%%%%%

Writing down the high order multiple scattering analogs of (\ref{Sscat}) in 
component form is very cumbersome. To simplify the notation, we introduce 
a matrix notation.
To this end, we assemble the basis states (\ref{P0rep})--(\ref{H0rep}) 
from {\sc Box} {\bf 4.3} into the $8 
\times 14$
array
\begin{equation} \label{Bab0aug}
|\re,k_y,{\bf 0}+)\equiv \left[ \begin{array}{ccc} |(\re,k_y,P),{\bf 0+}) & 
|(\re,k_y,S),{\bf 0+}) & |(\re,k_y,H),{\bf 0+}) \end{array}\right]
\end{equation}
%which is an augmented version of $6 \times 14$ array $|\re,k_y,{\bf 0+}\rangle$ from (\ref{Bab0Def}).
Also, we introduce inter-scatterer propagator ${{\bmi  H}}({\bmi  
R}_{i},{\bmi  R}_{j})$ defined as
\begin{equation} \label{Pdef1}
{{\bmi  H}}({\bmi  R}_{i},{\bmi  R}_{j})=
\left[{\begin{array}{ccc}{\bmi H}_{P}({\bmi R}_{i};{\bmi R}_{j})&{\bf 0}&{\bf 0}\\
{\bf 0}&{\bmi H}_{S}({\bmi R}_{i};{\bmi R}_{j})&{\bf 0}\\
{\bf 0}&{\bf 0}&{\bmi H}_{H}({\bmi R}_{i};{\bmi R}_{j})\end{array}}\right]
\end{equation}
with partitions
\begin{eqnarray} \label{PaDef0}
\lefteqn{{\bmi H}_{c }({\bmi R}_{i};{\bmi R}_{j})=\frac{1}{4 \rf{\rho} \omega^2}\times}\\
&&\!\!\!\!\!\!\!\!\!\!\!\!\!\!
\left[{\begin{array}{ccccc}
\scriptstyle 
{H}^{(1)}_{0}(\hat{K}_{c}{R}_{ij})&\scriptstyle\!\!\! -{H}^{(1)}_{1}(\hat{K}_{c}{R}_{ij})
Y_{1}(\hat{\bf R}_{ji})&\scriptstyle{H}^{(1)}_{2}(\hat{K}_{c}{R}_{ij})
Y_{2}(\hat{\bf R}_{ji})&\scriptstyle\!\!\! -{H}^{(1)}_{3}(\hat{K}_{c}{R}_{ij})
Y_{3}(\hat{\bf R}_{ji})&\scriptstyle{H}^{(1)}_{4}(\hat{K}_{c}{R}_{ij})
Y_{4}(\hat{\bf R}_{ji})\\ \scriptstyle {H}^{(1)}_{1}(\hat{K}_{c}{R}_{ij})
Y_{1}^{\ast}(\hat{\bf R}_{ji})&\scriptstyle{H}^{(1)}_{0}(\hat{K}_{c}{R}_{ij})&\scriptstyle\!\!\! -{H}^{(1)}_{1}
(\hat{K}_{c}{R}_{ij})Y_{1}(\hat{\bf R}_{ji})&\scriptstyle{H}^{(1)}_{2}(\hat{K}_{c}{R}_{ij})
Y_{2}(\hat{\bf R}_{ji})&\scriptstyle\!\!\! -{H}^{(1)}_{3}(\hat{K}_{c}{R}_{ij})
Y_{3}(\hat{\bf R}_{ji})\\ \scriptstyle {H}^{(1)}_{2}(\hat{K}_{c}{R}_{ij})
Y_{2}^{\ast}(\hat{\bf R}_{ji})&\scriptstyle{H}^{(1)}_{1}(\hat{K}_{c}{R}_{ij})
Y_{1}^{\ast}(\hat{\bf R}_{ji})&\scriptstyle{H}^{(1)}_{0}(\hat{K}_{c}{R}_{ij})&\scriptstyle\!\!\! -{H}^{(1)}_{1}
(\hat{K}_{c}{R}_{ij})
Y_{1}(\hat{\bf R}_{ji})&\scriptstyle{H}^{(1)}_{2}(\hat{K}_{c}{R}_{ij})
Y_{2}(\hat{\bf R}_{ji})\\ \scriptstyle {H}^{(1)}_{3}(\hat{K}_{c}{R}_{ij})
Y_{3}^{\ast}(\hat{\bf R}_{ji})&\scriptstyle{H}^{(1)}_{2}(\hat{K}_{c}{R}_{ij})
Y_{2}^{\ast}(\hat{\bf R}_{ji})&\scriptstyle{H}^{(1)}_{1}(\hat{K}_{c}{R}_{ij})
Y_{1}^{\ast}(\hat{\bf R}_{ji})&\scriptstyle{H}^{(1)}_{0}(\hat{K}_{c}{R}_{ij})&\scriptstyle\!\!\! -{H}^{(1)}_{1}
(\hat{K}_{c}{R}_{ij})Y_{1}(\hat{\bf R}_{ji})\\\!\! \scriptstyle {H}^{(1)}_{4}
(\hat{K}_{c}{R}_{ij})Y_{4}^{\ast}(\hat{\bf R}_{ji})&\scriptstyle{H}^{(1)}_{3}
(\hat{K}_{c}{R}_{ij})Y_{3}^{\ast}(\hat{\bf R}_{ji})&\scriptstyle{H}^{(1)}_{2}
(\hat{K}_{c}{R}_{ij})Y_{2}^{\ast}(\hat{\bf R}_{ji})&\scriptstyle{H}^{(1)}_{1}
(\hat{K}_{c}{R}_{ij})Y_{1}^{\ast}(\hat{\bf R}_{ji})&\scriptstyle{H}^{(1)}_{0}
(\hat{K}_{c}{R}_{ij})\end{array}}\!\!\right];\nonumber 
\end{eqnarray} 
for $c\in \{ P,H\}$, and
\begin{eqnarray} \label{PbDef0} 
\lefteqn{
{\bmi H}_{S}({\bmi R}_{i};{\bmi R}_{j})=\frac{1}{4 \rf{\rho} \omega^2}\times}\\
 && \left[{\begin{array}{cccc} 
\scriptstyle{H}^{(1)}_{0}(\hat{K}_{S}{R}_{ij})&\scriptstyle-{H}^{(1)}_{1}
(\hat{K}_{S}{R}_{ij})Y_{1}(\hat{\bf R}_{ji})&\scriptstyle-
{H}^{(1)}_{3}(\hat{K}_{S}{R}_{ij})Y_{3}(\hat{\bf R}_{ji})&
\scriptstyle{H}^{(1)}_{4}(\hat{K}_{S}{R}_{ij})Y_{4}(\hat{\bf R}_{ji})\\ 
\scriptstyle {H}^{(1)}_{1}(\hat{K}_{S}{R}_{ij})Y_{1}^{\ast}(\hat{\bf R}_{ji})&
\scriptstyle{H}^{(1)}_{0}(\hat{K}_{S}{R}_{ij})&\scriptstyle
{H}^{(1)}_{2}(\hat{K}_{S}{R}_{ij})Y_{2}(\hat{\bf R}_{ji})&\scriptstyle-
{H}^{(1)}_{3}(\hat{K}_{S}{R}_{ij})Y_{3}(\hat{\bf R}_{ji})\\ \scriptstyle 
{H}^{(1)}_{3}(\hat{K}_{S}{R}_{ij})Y_{3}^{\ast}(\hat{\bf R}_{ji})&
\scriptstyle{H}^{(1)}_{2}(\hat{K}_{S}{R}_{ij})Y_{2}^{\ast}(\hat{\bf R}_{ji})&\scriptstyle{H}^{(1)}_{0}(\hat{K}_{S}{R}_{ij})&\scriptstyle-{H}^{(1)}_{1}
(\hat{K}_{S}{R}_{ij})Y_{1}(\hat{\bf R}_{ji})\\ \scriptstyle 
{H}^{(1)}_{4}(\hat{K}_{S}{R}_{ij})Y_{4}^{\ast}(\hat{\bf R}_{ji})&
\scriptstyle{H}^{(1)}_{3}(\hat{K}_{S}{R}_{ij})Y_{3}^{\ast}(\hat{\bf R}_{ji})&
\scriptstyle{H}^{(1)}_{1}(\hat{K}_{S}{R}_{ij})Y_{1}^{\ast}(\hat{\bf R}_{ji})&\scriptstyle{H}^{(1)}_{0}(\hat{K}_{S}{R}_{ij}) 
\end{array}}\right].\nonumber 
\end{eqnarray} 
When $|k_y| > |\omega/c|$, we take the $+ve$ imaginary 
branch of the square root of $\hat{K}_c$ in (\ref{Kdef}), so that the Hankel 
functions  ${H}^{(1)}_{m}(\hat{K}_{c}r)$ in (\ref{PaDef0})--(\ref{PbDef0})
decay exponentially for large values of $r$.





Using (\ref{Bab0aug}) and (\ref{Pdef1}), we define the further arrays
\begin{equation}	\label{G1mat}
{\rf{\bmi  G}}({\bmi  x},{\bmi  x}^{\prime};k_y)
		\equiv -i\,
	|\re,k_y,{\bmi r}_{i}\rangle {{\bmi  H}}({\bmi  R}_{i};{\bmi  R}_{j})
	\langle\re,k_y,{\bmi r}_{j}|{\bmi J}_{6}
\end{equation}
which is the matrix analogue of (\ref{G1Trep}),
\begin{equation}	\label{Grmat}
{\rf{\bmi  G}_{r}}({\bmi  R}_{i},{\bmi  x}^{\prime};k_y)
		\equiv -i\,
	|\re,k_y,{\bf 0}+) {{\bmi  H}}({\bmi  R}_{i};{\bmi  R}_{j})
	\langle\re,k_y,{\bmi r}_{j}|{\bmi J}_{6},
\end{equation}
and
\begin{equation}	\label{Psmat}
{\rf{\bmi  P}_{s}}({\bmi  R}_{i},{\bmi  R}_{j};k_y)
		\equiv -i\,
	|\re,k_y,{\bmi r}_{i}\rangle {{\bmi  H}}({\bmi  R}_{i};{\bmi  R}_{j})
	(\re,k_y,{\bf 0}+|.
\end{equation}



The components (\ref{Sscat1}) of the single scattering interaction, may be 
assembled into an equivalent matrix form
\begin{eqnarray} \label{Sscat1mat}
\lefteqn{{\rf{\bmi G}}{\mit \Delta} {\bmi A} {\rf{\bmi G}}= {\rf{\bmi 
P}_{s}}{\mit \Delta} {\bmi C} {\rf{\bmi G}_{r}}}\nonumber\\
&&=\left(-i\,|\re,k_y,{\bmi r}_R\rangle {\bmi  H}({\bmi  R}_{R};{\bmi  R}_{j})
	\tilde{\bmi \Pi}|\re,-k_y,{\bf 0}+)^{T}\right)\nonumber\\
&&\times\left({\mit \Delta}{\bmi C}({\bmi R}_j;k_y)\,{\mit \Delta}^2  \right)
\left( -i\,|\re,k_y,{\bf 0}+) {\bmi  H}({\bmi  R}_{j};{\bmi  R}_{S})
	\tilde{\bmi \Pi}|\re,-k_y,{\bmi r}_S\rangle^{T}\right){\bmi 
	J}_{6}
\end{eqnarray}












%%%%%%%%%%%%%%%%%%%%%%%%%%%%%%%%%%%%%%%%%%%%%%%%%%%%%%%%%%%%%%%%%%%%%%%%%%%%%%%%%%%%%%%
%%%%%%%%%%%%%%%%%%%%%%%%%%%%%%%%%%%%%%%%%%%%%%%%%%%%%%%%%%%%%%%%%%%%%%%%%%%%%%%%%%%%%%%
%%%%%%%%%%%%%%%%%%%%%%%%%%%%%%%%%%%%%%%%%%%%%%%%%%%%%%%%%%%%%%%%%%%%%%%%%%%%%%%%%%%%%%%
%%%%%%%%%%%%%%%%%%%%%%%%%%%%%%%%%%%%%%%%%%%%%%%%%%%%%%%%%%%%%%%%%%%%%%%%%%%%%%%%%%%%%%%
\subsection{Double scattering}
Let us now consider the {\em double scattering interaction\/} 
\begin{eqnarray} \label{Dscat}
\lefteqn{{\rf{\bmi G}}{\mit \Delta} {\bmi A} {\rf{\bmi G}}
{\mit \Delta} {\bmi A}{\rf{\bmi G}}\equiv}\\
&&
\int_{{\mit \Omega}_j} \int_{{\mit \Omega}_k}{\rf{\bmi G}}({\bmi x}_{R};{\bmi x};k_y)
{\mit \Delta} {\bmi A}({\bmi
x},\partial_{x};k_y) {\rf{\bmi G}}({\bmi x},{\bmi x}^{\prime};k_y)
{\mit \Delta} {\bmi A}({\bmi
x}^{\prime},\partial_{x};k_y) {\rf{\bmi G}}({\bmi x}^{\prime},{\bmi 
x}_{S};k_y) 
{\rm d}{\bmi x}\, {\rm d}{\bmi x}^{\prime}\nonumber
\end{eqnarray}
between two small, 
square subregions of heterogeneity ${\mit \Omega}_k$ with centre at ${\bmi 
R}_k$ and  ${\mit \Omega}_j$ with centre at ${\bmi R}_j$. 
Using (\ref{G1Trep}) for the components of ${\rf{\bmi G}}$ in (\ref{Dscat}) 
gives
\begin{eqnarray} \label{Dscat1}
\lefteqn{{\rf{\bmi G}}{\mit \Delta} {\bmi A} {\rf{\bmi G}}
{\mit \Delta} {\bmi A}{\rf{\bmi G}}=
\sum\limits_{\stackrel{m,n}{c\in\{P,S,H\}}} 
\sum\limits_{\stackrel{m^{\prime},n^{\prime}}{c^{\prime}\in\{P,S,H\}}}
\sum\limits_{\stackrel{m^{\prime\prime},n^{\prime\prime}}
{c^{\prime\prime}\in\{P,S,H\}}}
\frac{i }{64\,\rf{\rho}\!\!{}^3\,\omega^6 }
 } \nonumber\\
&&\times
 |(\re,k_y,c,m),{\bmi r}_{R}\rangle {H}_{m-n}^{(1)}(\hat{ K}_{c}{R}_{Rj})
{Y}_{m-n}^{{}^\ast}(\hat{{\bmi R}}_{jR}) \nonumber\\
&&\times 
\int_{{\mit \Omega}_j}  \langle(\re,k_y,c,n),{\bmi x}-{\bmi R}_{j}|{\bmi J}_{6}
{\mit \Delta} {\bmi A}({\bmi
x},\partial_{x};k_y)  |(\re,k_y,c^{\prime},m^{\prime}),{\bmi x}-{\bmi R}_{j}\rangle 
{\rm d}{\bmi x}\nonumber\\
&&
\times {H}_{m^{\prime}-n^{\prime}}^{(1)}(\hat{ K}_{c^{\prime}}{R}_{jk})
{Y}_{m^{\prime}-n^{\prime}}^{{}^\ast}(\hat{{\bmi R}}_{kj}) \nonumber\\
  &&
\times \int_{{\mit \Omega}_k}  \langle(\re,k_y,c,n^{\prime}),{\bmi x}^{\prime}-{\bmi 
R}_{k}|{\bmi J}_{6}
{\mit \Delta} {\bmi A}({\bmi x}^{\prime},\partial_{x};k_y)  
|(\re,k_y,c^{\prime\prime},m^{\prime\prime}),{\bmi x}^{\prime}-{\bmi R}_{k}\rangle 
{\rm d}{\bmi x}^{\prime}\nonumber\\
&&
\times {H}_{m^{\prime\prime}-n^{\prime\prime}}^{(1)}(\hat{ 
K}_{c^{\prime\prime}}{R}_{kS})
{Y}_{m^{\prime\prime}-n^{\prime\prime}}^{{}^\ast}(\hat{{\bmi R}}_{Sk}) 
 \langle(\re,k_y,c^{\prime\prime},n^{\prime\prime}),{\bmi r}_{S}|{\bmi J}_{6}
\end{eqnarray} 
where we interpret the coupling integrals using the weak form 
of the Rayleigh approximation.


Once again, the components 
(\ref{Dscat1}) may be assembled into an equivalent matrix equation, which 
may be written as
\begin{eqnarray} \label{Dscat1mat}
 \lefteqn{{\rf{\bmi G}}{\mit \Delta} {\bmi A} {\rf{\bmi G}}
{\mit \Delta} {\bmi A}{\rf{\bmi G}}= {\rf{\bmi P}_{s}}{\mit 
\Delta} {\bmi C} {\rf{\bmi P}_{r,s}}
{\mit \Delta} {\bmi C}{\rf{\bmi G}_{r}}}\nonumber\\
&&=
\left(-i\,|\re,k_y,{\bmi r}_R\rangle {\bmi  H}({\bmi  R}_{R};{\bmi  R}_{j})
	(\re,-k_y,{\bf 0}+|\right)\left({\mit \Delta}{\bmi 
	C}({\bmi R}_j;k_y)\,{\mit \Delta}^2  \right)\nonumber\\
&&\!\!
\times
{\rf{\bmi  P}_{r,s}}({\bmi  R}_{j},{\bmi  R}_{k})
\left({\mit \Delta}{\bmi  C}({\bmi R}_j;k_y)\,{\mit \Delta}^2  \right)
\left( -i\,|\re,k_y,{\bf 0}+) {\bmi  H}({\bmi  R}_{i};{\bmi  R}_{j})
	\langle\re,k_y,{\bmi r}_S|\right){\bmi J}_{6}
\end{eqnarray}
to within the {\em weak Rayleigh approximation\/},
where the $8 \times 8 $ augmented coupling matrix  was given by 
(\ref{CdefAug}) and the $8 \times 8$ augmented propagator array
 ${\rf{\bmi  P}_{r,s}}$ has been written 
in the  matrix form
\begin{eqnarray}	\label{Prsmat}
	{\rf{\bmi  P}_{r,s}}({\bmi  R}_{i},{\bmi  R}_{j})
		\equiv	-i\,|\re,k_y,{\bf 0}+) {\bmi  H}({\bmi  R}_{i};{\bmi  R}_{j})
	(\re,k_y,{\bf 0}+|,
\end{eqnarray} 
and ${{\bmi  H}}({\bmi  R}_{i},{\bmi  R}_{j})$ was given by 
(\ref{Pdef1}).


\subsection{Evaluating the near field singularity}




We now consider the limiting case of (\ref{Dscat1mat}) where ${\mit 
\Omega}_j={\mit \Omega}_k={\mit \Omega}_{\varepsilon}$ is an 
{\em infinitesimal square\/} at the coordinate origin Fig.\ref{Omegaeps}. 
In this limiting case, the modified 
Green's tensor of (\ref{Prsmat}) is strongly singular, and we will 
have to be very careful how we interpret integrals involving 
(\ref{Prsmat}). We employ the small argument asymptotic approximations 
in Table \ref{HTable}, 
to the Hankel function entries in the diagonal blocks of 
(\ref{Pdef1}) that were defined in (\ref{PaDef0})--(\ref{PbDef0}).
 After a substantial amount of algebra, the leading order 
behaviour of the  Green's tensor 
in (\ref{Prsmat}) reduces to ${\tilde{\bmi  P}_{r,s}}({\bf 
0},{\bmi  r})$ given in {\sc Box} $\bf 4.4$. The $\delta-$function terms in 
${\tilde{\bmi  P}_{r,s}}({\bf 
0},{\bmi  r})$ are correction terms that arose when we took the partial 
differential operators ${\bmi L}_c, c\in \{P,S,H\}$ outside the integrals in (\ref{c4grepc}). We 
ignored these terms in the earlier representations of the Green's 
tensor because it was always assumed that the source and field point of 
the Green's tensor were distinct.  

%%%%%%%%%%%%%%%%%%%%%%%%%%%%%%%%%%%%%%%%%%%%%%%%%%%%%%%%%%%%%%%%%%%%%%%%%%%%%%
%%%%%%%%%%%%%%%%%%%%%%%%%%%%%%%%%%%%%%%%%%%%%%%%%%%%%%%%%%%%%%%%%%%%%%%%%%%%%%
%%%%%%%%%%%%%%%%%%%%%%%%%%%%%%%%%%%%%%%%%%%%%%%%%%%%%%%%%%%%%%%%%%%%%%%%%%%%%
\begin{figure}
\HideDisplacementBoxes
%\ShowDisplacementBoxes
\centerline{\BoxedEPSF{Omegaeps.EPSF scaled 800}}
\caption{The infinitesimal square $\Omega$ with an inscribed disc 
$D_{\varepsilon}$}
\protect\label{Omegaeps}
\end{figure}
%%%%%%%%%%%%%%%%%%%%%%%%%%%%%%%%%%%%%%%%%%%%%%%%%%%%%%%%%%%%%%%%%%%%%%%%%%%%%
%%%%%%%%%%%%%%%%%%%%%%%%%%%%%%%%%%%%%%%%%%%%%%%%%%%%%%%%%%%%%%%%%%%%%%%%%%%%%%%%%

\begin{table}
\begin{minipage} {90.0mm}
\begin{tabular*}{85mm}{ccl}\hline
$H^{(1)}_m(r)$&&\mbox{leading order asymptotics}\\ \hline
$H^{(1)}_0(r)$&$\approx$&$1 + i\,{{2}\over {\pi }}\,\left( \gamma_{E} + 
     {{ \log ({1\over 2}) + \log (r)  }} \right)+O(r)$\\
$H^{(1)}_1(r)$&$\approx$&$-{{2\,i}\over {\pi \,r}}+O(r)$\\
$H^{(1)}_2(r)$&$\approx$&$-{i\over {\pi }} - {4\,i\over {\pi \,{r^2}}}
+O(r)$\\
$H^{(1)}_3(r)$&$\approx$&$ - {2\,i\over {\pi \,r}} -{{16\,i}\over {\pi \,{r^3}}}
+O(r)$\\
$H^{(1)}_4(r)$&$\approx$&$-{{i}\over {2\,\pi }} - {8\,i\over {\pi 
\,{r^2}}}  - {{96\,i}\over {\pi \,{r^4}}}
+O(r)$\\
\hline
\end{tabular*}
\caption{\label{HTable} {$\gamma_{E}=0.5772156649015328606...$ {\em is the
 Euler-Mascheroni constant\/} (e.g. {\em Arfken\/}~(1985), p.284)}}
\end{minipage}
\end{table} 


The integral
\[
\lim\limits_{{ \mit \Omega }_{ \varepsilon }\ \rightarrow  \ 0}\ 
\int_{{\mit \Omega }_{ \varepsilon }}{\rf{\bmi  P}_{r,s}}({\bf 0},{\bmi  r}) \cdot 
{\mit \Delta}{\bmi C}({\bmi  r})\ 
{\rm d}{\bmi r}
\]
has a {\em strong singularity\/} because the Green's tensor 
${\rf{\bmi  P}_{r,s}}({\bf 0},{\bmi  r})$ is  $O(\|{\bmi  
r}\|^{-2})$ in the near field. To handle this integral
we employ the classical technique of {\em subtracting the singularity\/}, 
i.e.
\begin{eqnarray} \label{subint}
\lefteqn{\lim\limits_{{\mit  \Omega }_{ \varepsilon }\ \rightarrow  \ 0}\ 
\int_{{\mit  \Omega }_{ \varepsilon }}{\rf{\bmi  P}_{r,s}}({\bf 0},{\bmi  r}) \cdot 
{\mit \Delta}{\bmi C}({\bmi  r})\ 
{\rm d}{\bmi r}
=\lim\limits_{{\mit  \Omega }_{ \varepsilon }\ \rightarrow  \ 
0}\int_{{ \mit \Omega }_{ \varepsilon 
}}\left({{\rf{\bmi  P}_{r,s}}({\bf 0},{\bmi  r})-
{\tilde{\bmi  P}_{r,s}}({\bf 0},{\bmi  r})}\right) \cdot {\mit \Delta}{\bmi C}({\bmi  r})\ 
{\rm d}{\bmi r}}\nonumber\\
&&\qquad\qquad\qquad
+
\left(\lim\limits_{{ \mit \Omega }_{ \varepsilon }\ \rightarrow  \ 
0}\int_{{\mit  \Omega }_{ \varepsilon }}^{}{\tilde{\bmi  P}_{r,s}}({\bf 
0},{\bmi  r}){\rm d}{\bmi r}\right) \cdot{\mit \Delta}{\bmi C}({\bf 0}+)
\end{eqnarray}
as the domain of integration  ${ \mit \Omega }_{ \varepsilon }$ shrinks down 
to a point,
where the dominant near field term ${\tilde{\bmi  P}_{r,s}}({\bf 
0},{\bmi  r})$ of the Green's tensor is given in {\sc Box} {\bf 4.4}.
If the coupling matrix ${\mit \Delta}{\bmi C}({\bmi  r})$ is {\em 
Holder continuous\/} (i.e. it is smoother than just continuous , but not 
necessarily so smooth as to be differentiable)  in an open neighbourhood of the point ${\bmi r}={\bf 
0}$, then 
\[
\lim\limits_{{ \mit \Omega }_{ \varepsilon }\ \rightarrow  \ 
0}\int_{{\mit \Omega }_{ \varepsilon 
}}\left({{\rf{\bmi  P}_{r,s}}({\bf 0},{\bmi  r})-
{\tilde{\bmi  P}_{r,s}}({\bf 0},{\bmi  r})}\right) \cdot {\mit \Delta}{\bmi C}({\bmi  r})\ 
{\rm d}{\bmi r} \ \rightarrow  \ {\bf 0}_8
\]
in (\ref{subint}) so that we are left with
\begin{equation} \label{subint1}
\lim\limits_{{\mit \Omega }_{ \varepsilon }\ \rightarrow  \ 0}\ 
\int_{{\mit \Omega }_{ \varepsilon }}{\rf{\bmi  P}_{r,s}}({\bf 0},{\bmi  r}) \cdot 
{\mit \Delta}{\bmi C}({\bmi  r})\ 
{\rm d}{\bmi r}
=\left(\lim\limits_{{\mit \Omega }_{ \varepsilon }\ \rightarrow  \ 
0}\int_{{\mit \Omega }_{ \varepsilon }}^{}{\tilde{\bmi  P}_{r,s}}({\bf 
0},{\bmi  r}){\rm d}{\bmi r}\right) \cdot{\mit \Delta}{\bmi C}({\bf 0}+)
\end{equation}

%%%%%%%%%%%%%%%%%%%%%%%%%%%%%%%%%%%%%%%%%%%%%%%%%%%%%%%%%%%%%%%%%%%%%%%%%%%%%%%%%%%%%%%%%%%
%%%%%%%%%%%%%%%%%%%%%%%%%%%%%%%%%%%%%%%%%%%%%%%%%%%%%%%%%%%%%%%%%%%%%%%%%%%%%%%%%%%%%%%%%%%
%%%%%%%%%%%%%%%%%%%%%%%%%%%%%%%%%%%%%%%%%%%%%%%%%%%%%%%%%%%%%%%%%%%%%%%%%%%%%%%%%%%%%%%%%%%%%%%%%%%
\begin{figure}
\fbox{
\begin{minipage} {142.0mm}
{\sc Box} {\bf 4.4}  {\sc The leading order terms of the modified Green's tensor (\ref{Prsmat}) 
in the near field} 
 
\vspace{3mm}
\[
{\tilde{\bmi  P}_{r,s}}({\bf 0},{\bmi  r}) =
\]  
\[
-\delta(r)
   \left[\begin{array}{cccccccc} 
   0&0&0&0&0&0&0&0\\ 0&0&0&0&0&0&0&0\\ 0&0&0&0&0&0&0&0\\
   0&0&0&\rf{\lambda}+2 \rf{\mu}&0&0&0&0\\ 0&0&0&0&\rf{\mu}&0&0&0\\
   0&0&0&0&0&\rf{\mu}&0&0\\ 0&0&0&0&0&0&0&0\\
   0&0&0&0&0&0&0&0\end{array}\right]
+
{{\left( {\rf{\alpha}\!{}^2} - {\rf{\beta}\!{}^2} \right) \,\cos 4\,\theta }\over 
   {{\rf{\alpha}\!{}^2}\,\pi \,{r^2}}}
   \left[\begin{array}{cccccccc} 0&0&0&0&0&0&0&0\\ 0&0&0&0&0&0&0&0\\ 0&0&0&0&0&0&0&0\\
   0&0&0&\rf{\mu}&0&0&-{1\over 2}&0\\ 0&0&0&0&-\rf{\mu}&0&0&0\\
   0&0&0&0&0&0&0&0\\ 0&0&0&-{1\over 2}&0&0&{1\over {4\,\rf{\mu}}}&0\\
   0&0&0&0&0&0&0&0\end{array}\right]
\]
  

   
\[
+
{{\left( {\rf{\alpha}\!{}^2} - {\rf{\beta}\!{}^2} \right) \,\sin 4\,\theta }\over 
   {{\rf{\alpha}\!{}^2}\,\pi \,{r^2}}}
\left[\begin{array}{cccccccc} 0&0&0&0&0&0&0&0\\ 0&0&0&0&0&0&0&0\\ 0&0&0&0&0&0&0&0\\
   0&0&0&0&\rf{\mu}&0&0&0\\ 0&0&0&\rf{\mu}&0&0&-{1\over 2}&0\\
   0&0&0&0&0&0&0&0\\ 0&0&0&0&-{1\over 2}&0&0&0\\ 0&0&0&0&0&0&0&0\end{array}\right]
\]
  

\[
+
{{\cos 2\,\theta }\over {{r^2}}}
\left[\begin{array}{cccccccc} 0&0&0&0&0&0&0&0\\ 0&0&0&0&0&0&0&0\\ 0&0&0&0&0&0&0&0\\
   0&0&0&{{2\,{\rf{\beta}\!{}^2}\,\left( {\rf{\alpha}\!{}^2} - {\rf{\beta}\!{}^2} \right) \,
       {\rf{\rho}}}\over {{\rf{\alpha}\!{}^2}\,\pi }}&0&0&
   {{-{\rf{\alpha}\!{}^2} + 2\,{\rf{\beta}\!{}^2}}\over {2\,{\rf{\alpha}\!{}^2}\,\pi }}&0\\
   0&0&0&0&0&0&0&0\\ 0&0&0&0&0&
   {{{\rf{\beta}\!{}^2}\,{\rf{\rho}}}\over {2\,\pi }}&0&0\\
   0&0&0&{{-{\rf{\alpha}\!{}^2} + 2\,{\rf{\beta}\!{}^2}}\over 
     {2\,{\rf{\alpha}\!{}^2}\,\pi }}&0&0&
   {{-1}\over {2\,{\rf{\alpha}\!{}^2}\,\pi \,{\rf{\rho}}}}&0\\
   0&0&0&0&0&0&0&{{-1}\over {2\,{\rf{\beta}\!{}^2}\,\pi \,{\rf{\rho}}}}\end{array}\right]
\]
  
  

\[
+
{{\sin 2\,\theta }\over {{r^2}}}  
\left[\begin{array}{cccccccc} 
0&0&0&0&0&0&0&0\\ 0&0&0&0&0&0&0&0\\ 0&0&0&0&0&0&0&0\\
   0&0&0&0&{{{\rf{\beta}\!{}^2}\,\left( {\rf{\alpha}\!{}^2} - {\rf{\beta}\!{}^2} \right) \,
       {\rf{\rho}}}\over {{\rf{\alpha}\!{}^2}\,\pi }}&0&0&0\\
   0&0&0&0&0&0&{{{\rf{\beta}\!{}^2}}\over {2\,{\rf{\alpha}\!{}^2}\,\pi }}&0\\
   0&0&0&0&0&0&0&{1\over {2\,\pi }}\\
   0&0&0&0&{{{\rf{\beta}\!{}^2}}\over {2\,{\rf{\alpha}\!{}^2}\,\pi }}&0&0&0\\
   0&0&0&0&0&{1\over {2\,\pi }}&0&0\end{array}\right]
\] 
\vspace{3mm}
\end{minipage}
}
\end{figure}
%%%%%%%%%%%%%%%%%%%%%%%%%%%%%%%%%%%%%%%%%%%%%%%%%%%%%%%%%%%%%%%%%%%%%%%%%%%%%%%%%%%%%%%%%%%
%%%%%%%%%%%%%%%%%%%%%%%%%%%%%%%%%%%%%%%%%%%%%%%%%%%%%%%%%%%%%%%%%%%%%%%%%%%%%%%%%%
%%%%%%%%%%%%%%%%%%%%%%%%%%%%%%%%%%%%%%%%%%%%%%%%%%%%%%%%%%%%%%%%%%%%%%%%%%%%%%%%%%%%%%%%%%%
%%%%%%%%%%%%%%%%%%%%%%%%%%%%%%%%%%%%%%%%%%%%%%%%%%%%%%%%%%%%%%%%%%%%%%%%%%%%%%%%%%%%%%%%%%%
%%%%%%%%%%%%%%%%%%%%%%%%%%%%%%%%%%%%%%%%%%%%%%%%%%%%%%%%%%%%%%%%%%%%%%%%%%%%%%%%%%%%%%%%%%%

We divide the infinitesimal square ${\mit \Omega }_{ \varepsilon }= 
\{(z,x): -\varepsilon \le z,x \le \varepsilon\} $ in (\ref{subint1}) into the sum 
of an inscribed {\em disc\/} ${\mit D}_{ \varepsilon }= 
\{(z,x): \sqrt{z^2 + x^2} \le \varepsilon\} $ and the 
remainder ${\mit \Omega}_{ \varepsilon }\backslash{\mit D}_{ \varepsilon 
}$ (Fig.\ref{Omegaeps}),
so that (\ref{subint1}) may be written as
\begin{eqnarray} \label{subint2}
\lefteqn{\lim\limits_{{\mit \Omega }_{ \varepsilon }\ \rightarrow  \ 0}\ 
\int_{{\mit \Omega }_{ \varepsilon }}{\rf{\bmi  P}_{r,s}}({\bf 0},{\bmi  r}) \cdot 
{\mit \Delta}{\bmi C}({\bmi  r})\ 
{\rm d}{\bmi r}
=}\\
&&\left(\lim\limits_{{\mit D }_{ \varepsilon }\ \rightarrow  \ 
0}\int_{{\mit D }_{ \varepsilon }}{\tilde{\bmi  P}_{r,s}}({\bf 
0},{\bmi  r}){\rm d}{\bmi r}\right) \cdot{\mit \Delta}{\bmi C}({\bf 
0}+)
+
\left(\lim\limits_{{\mit \Omega }_{ \varepsilon }\ \rightarrow  \ 
0}\int_{{\mit \Omega }_{ \varepsilon }\backslash
{\mit D}_{ \varepsilon }}{\tilde{\bmi  P}_{r,s}}({\bf 
0},{\bmi  r}){\rm d}{\bmi r}\right) \cdot{\mit \Delta}{\bmi C}({\bf 
0}+)\nonumber
\end{eqnarray}
Once again, we employ classical techniques to evaluate the two {\em 
strongly singular \/} integrals on the right hand side of (\ref{subint2}).
First, we consider the integral
\[
\lim\limits_{{\mit D }_{ \varepsilon }\ \rightarrow  \ 
0}\int_{{\mit D}_{ \varepsilon }}{\tilde{\bmi  P}_{r,s}}({\bf 
0},{\bmi  r}){\rm d}{\bmi r}
\]
The $\delta-$function contributions to this integral that were indicated 
in {\sc Box} {\bf 4.4} are trivial to integrate. To evaluate the
$\int_{{\mit \Omega}_{\varepsilon}}{{\cos m\,\theta }\over {{r^2}}}$ and 
$\int_{{\mit \Omega}_{\varepsilon}}{{\sin m\,\theta }\over {{r^2}}}, m\in 
\{2,4\}$ integrals,  we employ the 
$2-$D {\em divergence theorem\/} in cylindrical coordinates, i.e.
given the $2-$D vector field ${\bmi v}=({v}_{r},{v}_{\theta })$ in 
cylindrical coordinates, then
\begin{equation} \label{divthm}
\int_{{\mit D }_{ \varepsilon }}^{} \nabla \cdot   
{\bmi v} {\rm d}{\bmi r} \equiv \int_{  0}^{  2 \pi }  {\rm d} \theta 
\int_{  0}^{\varepsilon }  r\ 
{\rm d}r\left[{{\frac{1}{r}}{\frac{ \partial   \,{v}_{ \theta }}{ 
\partial   \, \theta }}+{\frac{1}{r}}{\frac{ \partial   
\,r{v}_{r}}{ \partial   \, r }}}\right] \equiv 
\left.\int_{  0}^{  2 \pi }r {  v}_{  
r}  {\rm d} \theta \right|_{r=\varepsilon}
\end{equation} 
 Using
\[
{\frac{\cos\ \left({m\rm \theta 
}\right)}{{r}^{2}}}={\frac{1}{r}}{\frac{\rm \partial \rm \mit \,}{\rm 
\partial \rm \mit \,\rm \theta }}{\frac{\sin\ \left({m\rm \theta 
}\right)}{m\ r}} \qquad \mbox{and} \qquad 
{\frac{\sin\ \left({m\rm \theta 
}\right)}{{r}^{2}}}=-{\frac{1}{r}}{\frac{\rm \partial \rm \mit \,}{\rm 
\partial \rm \mit \,\rm \theta }}{\frac{\cos\ \left({m\rm \theta 
}\right)}{m\ r}}
\] 
it follows from (\ref{divthm}) that
\[
\int_{{\mit D}_{\varepsilon}}{{\cos m\,\theta }\over {{r^2}}} {\rm 
d}{\bmi r}\equiv 0, \quad \mbox{and}
\quad \int_{{\mit D}_{\varepsilon}}{{\sin m\,\theta }\over {{r^2}}} {\rm 
d}{\bmi r}\equiv 0, \quad m\in\{2,4\}
\]
because the $v_r$ components in (\ref{divthm}) vanish. Consequently, we 
are left with the $\delta-$function terms, i.e.
\begin{equation} \label{intD}
\lim\limits_{{\mit D }_{ \varepsilon }\ \rightarrow  \ 
0}\int_{{\mit D}_{ \varepsilon }}{\tilde{\bmi  P}_{r,s}}({\bf 
0},{\bmi  r}){\rm d}{\bmi r}
\equiv
-   \left[\begin{array}{cccccccc} 
   0&0&0&0&0&0&0&0\\ 0&0&0&0&0&0&0&0\\ 0&0&0&0&0&0&0&0\\
   0&0&0&\rf{\lambda}+2 \rf{\mu}&0&0&0&0\\ 0&0&0&0&\rf{\mu}&0&0&0\\
   0&0&0&0&0&\rf{\mu}&0&0\\ 0&0&0&0&0&0&0&0\\
   0&0&0&0&0&0&0&0\end{array}\right].
\end{equation}


Next we consider the integral
\begin{equation} \label{intDm}
\lim\limits_{{\mit \Omega }_{ \varepsilon }\ \rightarrow  \ 
0}\int_{{\mit \Omega }_{ \varepsilon }\backslash{\mit D}_{ \varepsilon }}{\tilde{\bmi  P}_{r,s}}({\bf 
0},{\bmi  r}){\rm d}{\bmi r}
\end{equation}
from (\ref{subint2}). In this case the $\delta-$function terms in 
${\tilde{\bmi  P}_{r,s}}$ are 
excluded from the domain of integration, and we need only consider the 
integrals
\[
\int_{{\mit \Omega}_{\varepsilon}\backslash{\mit D}_{ \varepsilon }}{{\cos m\,\theta }\over {{r^2}}} {\rm 
d}{\bmi r}, \quad \mbox{and}
\quad \int_{{\mit \Omega}_{\varepsilon}\backslash{\mit D}_{ \varepsilon }}{{\sin m\,\theta }\over {{r^2}}} {\rm 
d}{\bmi r}, \quad m\in\{2,4\}.
\]
But $\int_{{\mit \Omega}_{\varepsilon}\backslash{\mit D}_{ \varepsilon }}$
 may be split into the sum of the {\em quadrant\/} integrals
\begin{eqnarray} \label{quadint}
\int_{{\mit \Omega}_{\varepsilon}\backslash{\mit D}_{ \varepsilon }}\left[ \cdot 
\right] {\rm d}{\bmi r}&\equiv&
\int_{- \pi   /4}^{ \pi   /4}{\rm d} \theta 
\int_{\varepsilon }^{\varepsilon   /\cos\ \left({ \theta 
}\right)}\left[ \cdot 
\right] r\ {\rm d}r\ + \
\int_{3 \pi   /4}^{5 \pi   
/4}{\rm d} \theta \int_{\varepsilon }^{  - \varepsilon   
/\cos\ \left({ \theta }\right)}\left[ \cdot 
\right] r\ {\rm d}r \nonumber\\
&&  +\ \int_{ \pi   
/4}^{3 \pi   /4}{\rm d} \theta \int_{\varepsilon }^{\varepsilon  
 /\sin\ \left({ \theta }\right)}\left[ \cdot 
\right] r\ {\rm d}r\ +\ \int_{5 \pi 
  /4}^{7 \pi   /4}{\rm d} \theta \int_{\varepsilon }^{ 
 - \varepsilon   /\sin\ \left({ \theta }\right)}\left[ \cdot 
\right] 
r\ {\rm d}r.
\end{eqnarray} 
Using the quadrant integrals (and integration by parts where 
necessary) we get
\[
\int_{{\mit \Omega}_{\varepsilon}\backslash{\mit D}_{ \varepsilon }}{{\cos 4\,\theta }\over {{r^2}}}  
{\rm d}{\bmi r}= {{\pi - 4} \over 2}
\]
and
\[
\int_{{\mit \Omega}_{\varepsilon}\backslash{\mit D}_{ \varepsilon }}{{\cos 2\,\theta }\over {{r^2}}}  {\rm d}{\bmi r}= 
\int_{{\mit \Omega}_{\varepsilon}\backslash{\mit D}_{ \varepsilon }}{{\sin 
4\,\theta }\over {{r^2}}}  {\rm d}{\bmi r}=
\int_{{\mit \Omega}_{\varepsilon}\backslash{\mit D}_{ \varepsilon }}{{\sin 
2\,\theta }\over {{r^2}}}  {\rm d}{\bmi r}=0.
\]
so that (\ref{intDm}) becomes
\begin{equation} \label{intDm1}
\lim\limits_{{\mit \Omega }_{ \varepsilon }\ \rightarrow  \ 
0}\int_{{\mit \Omega }_{ \varepsilon }\backslash{\mit D}_{ \varepsilon }}{\tilde{\bmi  P}_{r,s}}({\bf 
0},{\bmi  r}){\rm d}{\bmi r}=
{{({\pi - 4})( {\rf{\alpha}\!{}^2} - {\rf{\beta}\!{}^2})}\over 
   {2\, \pi\,{\rf{\alpha}\!{}^2}}}
   \left[\begin{array}{cccccccc} 0&0&0&0&0&0&0&0\\ 0&0&0&0&0&0&0&0\\ 0&0&0&0&0&0&0&0\\
   0&0&0&\rf{\mu}&0&0&-{1\over 2}&0\\ 0&0&0&0&-\rf{\mu}&0&0&0\\
   0&0&0&0&0&0&0&0\\ 0&0&0&-{1\over 2}&0&0&{1\over {4\,\rf{\mu}}}&0\\
   0&0&0&0&0&0&0&0\end{array}\right]
\end{equation}

Using (\ref{intD}) and (\ref{intDm1}) in the right hand side of 
(\ref{subint2}) recovers (\ref{subint1})
with
\begin{samepage}
\begin{eqnarray}  \label{Gepsint}
\lefteqn{
\lim\limits_{{\mit \Omega }_{ \varepsilon }\ \rightarrow  \ 
0}\int_{{\mit \Omega }_{ \varepsilon }}{\tilde{\bmi  P}_{r,s}}({\bf 
0},{\bmi  r}){\rm d}{\bmi r}=}\\
&&
-  \left[\begin{array}{cccccccc} 
   0&0&0&0&0&0&0&0\\ 0&0&0&0&0&0&0&0\\ 0&0&0&0&0&0&0&0\\
   0&0&0&\rf{\lambda}+2 \rf{\mu}&0&0&0&0\\ 0&0&0&0&\rf{\mu}&0&0&0\\
   0&0&0&0&0&\rf{\mu}&0&0\\ 0&0&0&0&0&0&0&0\\
   0&0&0&0&0&0&0&0\end{array}\right]
+
{{\left({\pi - 4}\right)( {\rf{\alpha}\!{}^2} - {\rf{\beta}\!{}^2} )}\over 
   {2\,{\rf{\alpha}\!{}^2}\,\pi}}
   \left[\begin{array}{cccccccc} 0&0&0&0&0&0&0&0\\ 0&0&0&0&0&0&0&0\\ 0&0&0&0&0&0&0&0\\
   0&0&0&\rf{\mu}&0&0&-{1\over 2}&0\\ 0&0&0&0&-\rf{\mu}&0&0&0\\
   0&0&0&0&0&0&0&0\\ 0&0&0&-{1\over 2}&0&0&{1\over {4\,\rf{\mu}}}&0\\
   0&0&0&0&0&0&0&0\end{array}\right]\nonumber
\end{eqnarray}
\end{samepage}

\subsection{Summing the multiple scattering series at a point}

We now return to the {\em double scattering integral\/} of (\ref{Dscat1mat}), 
and considering the limiting case where ${\mit 
\Omega}_k\rightarrow{\mit \Omega}_j$, so that the two small subregions 
coincide.  To within the {\em weak 
Rayleigh approximation\/} of (\ref{Dscat1mat}) , we have 
\begin{eqnarray} \label{Dscat2mat}
{\rf{\bmi G}}{\mit \Delta} {\bmi A} {\rf{\bmi G}}
{\mit \Delta} {\bmi A}{\rf{\bmi G}}&=&
\left(-i\,|\re,k_y,{\bf 0}+\rangle {\bmi  P}({\bmi  R}_{R};{\bmi  R}_{j})
	\tilde{\bmi \Pi}|\re,-k_y,{\bf 0}+)^{T}\right)\cdot\left({\mit \Delta}{\bmi 
	C}({\bmi R}_j;k_y)\,{\mit \Delta}^2  \right)
\nonumber\\
&&\times
\int_{{\mit \Omega }_{j}}{\tilde{\bmi  
G}}({\bmi 
R}_j,{\bmi  R}_j+{\bmi  r}_j){\rm d}{\bmi r}_j \cdot{\mit \Delta}{\bmi 
	C}({\bmi R}_j;k_y)\nonumber\\
&&\times
\left( -i\,|\re,k_y,{\bf 0}+) {\bmi  P}({\bmi  R}_{i};{\bmi  R}_{j})
	\tilde{\bmi \Pi}|\re,-k_y,{\bf 0}+\rangle^{T}\right){\bmi J}_{6}
\end{eqnarray}

Similarly, the {\em  order-$n$ multiple scattering interaction\/} 
\[
{\rf{\bmi G}}{\mit \Delta} {\bmi A} \left({\rf{\bmi G}}
{\mit \Delta} {\bmi A}\right)^n {\rf{\bmi G}}
\]
at 
${\mit \Omega}_j$, may be written as
\begin{eqnarray} \label{Dscatnmat}
\lefteqn{{\rf{\bmi G}}{\mit \Delta} {\bmi A} \left({\rf{\bmi G}}
{\mit \Delta} {\bmi A}\right)^n {\rf{\bmi G}}=
\left(-i\,|\re,k_y,{\bf 0}+\rangle {\bmi  P}({\bmi  R}_{R};{\bmi  R}_{j})
	\tilde{\bmi \Pi}|\re,-k_y,{\bf 0}+)^{T}\right)}\nonumber\\
&&	\times\left({\mit \Delta}{\bmi 
	C}({\bmi R}_j;k_y)\,{\mit \Delta}^2  \right)
\left(\int_{{\mit \Omega }_{j}}{\tilde{\bmi  
G}}({\bmi 
R}_j,{\bmi  R}_j+{\bmi  r}_j){\rm d}{\bmi r}_j \cdot{\mit \Delta}{\bmi 
	C}({\bmi R}_j;k_y)\right)^n\nonumber\\
&&\times
\left( -i\,|\re,k_y,{\bf 0}+) {\bmi  P}({\bmi  R}_{i};{\bmi  R}_{j})
	\tilde{\bmi \Pi}|\re,-k_y,{\bf 0}+\rangle^{T}\right){\bmi J}_{6}
\end{eqnarray}
If we sum all the multiple scattering interactions within ${ \mit \Omega }_{j}$  we get
\begin{samepage}
\begin{eqnarray} \label{msomegaj}
\lefteqn{
 \sum\limits_{ n\ =\ 0}^{ \infty } {\rf{\bmi G}}{\mit \Delta} {\bmi A} \left({\rf{\bmi G}}
{\mit \Delta} {\bmi A}\right)^n {\rf{\bmi G}}=
\left(-i\,|\re,k_y,{\bf 0}+\rangle {\bmi  P}({\bmi  R}_{R};{\bmi  R}_{j})
	\tilde{\bmi \Pi}|\re,-k_y,{\bf 0}+)^{T}\right)}
\nonumber\\
&&
\times \left({\mit \Delta}{\bmi 
	C}({\bmi R}_j;k_y)\,{\mit \Delta}^2  \right)
\sum\limits_{ n\ =\ 0}^{ \infty }\left(\int_{{\mit \Omega }_{j}}{\tilde{\bmi  
G}}({\bmi 
R}_j,{\bmi  R}_j+{\bmi  r}_j){\rm d}{\bmi r}_j \cdot{\mit \Delta}{\bmi 
	C}({\bmi R}_j;k_y)\right)^n\nonumber\\
&&\times
\left( -i\,|\re,k_y,{\bf 0}+) {\bmi  P}({\bmi  R}_{i};{\bmi  R}_{j})
	\tilde{\bmi \Pi}|\re,-k_y,{\bf 0}+\rangle^{T}\right){\bmi J}_{6}
\end{eqnarray}
\end{samepage}


%%%%%%%%%%%%%%%%%%%%%%%%%%%%%%%%%%%%%%%%%%%%%%%%%%%%%%%%%%%%%%%%%%%%%%%%%%%%%%%%%%%%%%%%%%%%%%%
%%%%%%%%%%%%%%%%%%%%%%%%%%%%%%%%%%%%%%%%%%%%%%%%%%%%%%%%%%%%%%%%%%%%%%%%%%%%%%%%%%%%%%%%%%%%%%%%
%%%%%%%%%%%%%%%%%%%%%%%%%%%%%%%%%%%%%%%%%%%%%%%%%%%%%%%%%%%%%%%%%%%%%%%%%%%%%%%%%%%%%%%%%%%%%%%

\begin{table}
\begin{tabular*}{136mm}{ccl}\hline\hline
&&\\

${{\mit \Delta} \tilde{\bmi C}}_{33}$&=&${\mit \Delta} \zeta\,{k_{y}^2}- 
{\mit \Delta} \rho\,{\omega^2}+{{k_{y}^2} \over {\it D}}\left(
       {{\mit \Delta}\chi^2}\,{\tilde{\zeta}} - 
       4\,{\mit \Delta}\chi\,{\mit \Delta} \gamma\,{\tilde{\zeta}}\,\rf{\mu} + 
       4\,{{\mit \Delta} \gamma^2}\,{\tilde{\zeta}}\,{\rf{\mu}\!{}^{2}}  
        \right.$\\
&&      $  \left. -{\rf{\alpha}\!{}^2}\,{\mit \Delta}a\,{{\mit \Delta}\chi^2}\,{\tilde{\zeta}}\,
        \rf{\rho} -2\,{\rf{\alpha}\!{}^2}\,{\mit \Delta}\chi\,{{\mit \Delta} \gamma^2}\,{\tilde{\zeta}}\,
        \rf{\rho} + 
       {\rf{\alpha}\!{}^2}\,{{\mit \Delta} \gamma^2}\,
       {\mit \Delta} \zeta\,{\tilde{\zeta}}\,
        \rf{\rho} - 4\,{\rf{\alpha}\!{}^2}\,{{\mit \Delta} \gamma^2}\,
        {\rf{\mu}\!{}^{2}}\,\rf{\rho}\right)$\\

&&\\
       
${{\mit \Delta} \tilde{\bmi C}}_{43}$&$=$& $- {i\,k_{y}\over {\it D}} 
{\left( -{\mit \Delta}\chi\,{\mit \Delta} \gamma\,{\tilde{\zeta}}   + 
         {\mit \Delta} \gamma\,{\mit \Delta} \zeta\,{\tilde{\zeta}} - 
         2\,{\mit \Delta}a\,{\mit \Delta}\chi\,{\tilde{\zeta}}\,\rf{\mu} - 
         2\,{{\mit \Delta} \gamma^2}\,{\tilde{\zeta}}\,\rf{\mu} - 4\,{\mit \Delta} \gamma\,{\rf{\mu}\!{}^{2}} \
         \right) }$\\

&&\\


${{\mit \Delta} \tilde{\bmi C}}_{73}$&$=$& $-{2\,i\,k_{y}\,\rf{\mu} \over  {\it D}}
\left(
 -{\mit \Delta}\chi\,{\mit \Delta} \gamma\,{\tilde{\zeta}} 
      + {\mit \Delta} \gamma\,{\mit \Delta} \zeta\,{\tilde{\zeta}} - 2\,{\mit \Delta}\chi\,\rf{\mu} 
      -  2\,{\mit \Delta}a\,{\mit \Delta}\chi\,{\tilde{\zeta}}\,\rf{\mu}  
      \right.$\\
&&        $\left.- 2\,{{\mit \Delta} \gamma^2}\,{\tilde{\zeta}}\,\rf{\mu} + 
         2\,{\rf{\alpha}\!{}^2}\,{\mit \Delta}a\,{\mit \Delta}\chi\,\rf{\mu}\,\rf{\rho} + 
         2\,{\rf{\alpha}\!{}^2}\,{{\mit \Delta} 
         \gamma^2}\,\rf{\mu}\,\rf{\rho} \right) $\\

&&\\

         
${{\mit \Delta} \tilde{\bmi C}}_{44}$&$=$& ${{{{\mit \Delta} \gamma^2}\,{\tilde{\zeta}}
 + {\mit \Delta}a\,{\mit \Delta} \zeta\,{\tilde{\zeta}} - 4\,{\mit 
\Delta}a\,{\rf{\mu}\!{}^{2}}}\over {{\it D}}} $\\

&&\\

${{\mit \Delta} \tilde{\bmi C}}_{74}$&$=$&${{2\,\rf{\mu}\,
       \left( {{\mit \Delta} \gamma^2}\,{\tilde{\zeta}} + {\mit \Delta}a\,{\mit \Delta} \zeta\,{\tilde{\zeta}} + 
         2\,{\mit \Delta} \gamma\,\rf{\mu} \right) }\over {{\it D}}}$\\

&&\\

${{\mit \Delta} \tilde{\bmi C}}_{55}$&$=$&${{{\mit \Delta}b}\over 
     {{\mit \Delta}b\,{\tilde{\zeta}} + \rf{\mu}\,{\mit \Delta}b -1 }}$\\

&&\\

${{\mit \Delta} \tilde{\bmi C}}_{66}$&$=$&${{{\mit \Delta}b}\over 
{\rf{\mu}\,{\mit \Delta}b  -1}}$\\

&&\\

${{\mit \Delta} \tilde{\bmi C}}_{77}$&$=$&${{4\,{\rf{\mu}\!{}^{2}}\,\left( {\mit \Delta} \zeta + {{\mit \Delta} \gamma^2}\,{\tilde{\zeta}} + 
         {\mit \Delta}a\,{\mit \Delta} \zeta\,{\tilde{\zeta}} - 
         {\rf{\alpha}\!{}^2}\,{{\mit \Delta} \gamma^2}\,\rf{\rho} - 
         {\rf{\alpha}\!{}^2}\,{\mit \Delta}a\,{\mit \Delta} \zeta\,\rf{\rho} \right) }\over 
         {{\it D}}}$\\

&&\\

$D$&$=$&$4\,{\mit \Delta} \gamma\,{\tilde{\zeta}}\,\rf{\mu} + 
  4\,{\rf{\mu}\!{}^{2}} + 4\,{\mit \Delta}a\,{\tilde{\zeta}}\,{\rf{\mu}\!{}^{2}} + 
  {\rf{\alpha}\!{}^2}\,{{\mit \Delta} 
  \gamma^2}\,{\tilde{\zeta}}\,\rf{\rho} +$\\
&& ${\rf{\alpha}\!{}^2}\,{\mit \Delta}a\,{\mit \Delta} \zeta\,{\tilde{\zeta}}\,\rf{\rho} - 
  4\,{\rf{\alpha}\!{}^2}\,{\mit \Delta}a\,{\rf{\mu}\!{}^{2}}\,\rf{\rho}
  -{\mit \Delta} \zeta\,{\tilde{\zeta}}  $\\

&&\\

$\tilde{\zeta}$&$=$&${{\left( \pi-4   \right) \,\rf{\zeta}}\over {8\,\pi 
}}$\\ 

&&\\ \hline\hline
\end{tabular*}
\caption{\label{CeffTable} 
{ The elements of the effective coupling 
matrix ${\mit \Delta}{\mbox{\mib C}_{eff}}$
}}
\end{table} 
%%%%%%%%%%%%%%%%%%%%%%%%%%%%%%%%%%%%%%%%%%%%%%%%%%%%%%%%%%%%%%%%%%%%%%%%%%%%%%%%%%%%%%%%%%%%%%%
%%%%%%%%%%%%%%%%%%%%%%%%%%%%%%%%%%%%%%%%%%%%%%%%%%%%%%%%%%%%%%%%%%%%%%%%%%%%%%%%%%%%%%%%%%%%%%%%
%%%%%%%%%%%%%%%%%%%%%%%%%%%%%%%%%%%%%%%%%%%%%%%%%%%%%%%%%%%%%%%%%%%%%%%%%%%%%%%%%%%%%%%%%%%%%%%


\subsection{The effective scattering parameters}
By comparing the multiple scattering expression (\ref{msomegaj}) with the 
single scattering expression (\ref{Sscat1mat}), we observe the correspondence
\[
{\mit \Delta}{\bmi C}({\bmi R}_j;k_y)
\Leftrightarrow
{\mit \Delta}{\bmi 
	C}({\bmi R}_j;k_y)
\sum\limits_{ n\ =\ 0}^{ \infty }\left(\int_{{\mit \Omega }_{j}}{\tilde{\bmi  
G}}({\bmi 
R}_j,{\bmi  R}_j+{\bmi  r}_j){\rm d}{\bmi r}_j \cdot{\mit \Delta}{\bmi 
	C}({\bmi R}_j;k_y)\right)^n
\]
We introduce an {\em effective coupling matrix\/} ${\mit \Delta}{\bmi 
C}_{eff}$ defined by
\begin{eqnarray} \label{dCeff0}
\lefteqn{{\mit \Delta}{\bmi 
	C}_{eff}({\bmi R}_j;k_y)\equiv}\nonumber\\
	&&
\lim\limits_{{ \mit \Omega }_{j} \rightarrow   0}\! 
{\mit \Delta}{\bmi 	C}({\bmi R}_j;k_y)\!
\sum\limits_{ n\ =\ 0}^{ \infty }\left(\int_{{\mit \Omega }_{j}}{\tilde{\bmi  
G}}({\bmi 
R}_j,{\bmi  R}_j+{\bmi  r}_j){\rm d}{\bmi r}_j \cdot{\mit \Delta}{\bmi 
	C}({\bmi R}_j;k_y)\right)^n
\end{eqnarray}
and using ${\mit \Delta}{\bmi	C}_{eff}({\bmi R}_j;k_y)$ 
in place of the coupling 
matrix ${\mit \Delta}{\bmi C}({\bmi R}_j;k_y)$ in (\ref{Sscat1mat}),  
the single scattering interaction
\begin{eqnarray} \label{Sscat2mat}
\lefteqn{{\rf{\bmi G}}{\mit \Delta} {\bmi A}_{eff} {\rf{\bmi G}}=
\left(-i\,|\re,k_y,{\bmi r}_R\rangle {\bmi  P}({\bmi  R}_{R};{\bmi  R}_{j})
	\tilde{\bmi \Pi}|\re,-k_y,{\bf 0}+)^{T}\right)}\nonumber\\
&&\times\left({\mit \Delta}\tilde{\bmi 
	C}({\bmi R}_j;k_y)\,{\mit \Delta}^2  \right)
\left( -i\,|\re,k_y,{\bf 0}+) {\bmi  P}({\bmi  R}_{j};{\bmi  R}_{S})
	\tilde{\bmi \Pi}|\re,-k_y,{\bmi r}_S\rangle^{T}\right){\bmi J}_{6}
\end{eqnarray}
where ${\mit \Delta} {\bmi A}_{eff}$ is ${\mit \Delta} {\bmi A}$ from (\ref{dAdef}) but 
constructed using the {\em effective parameter contrasts\/} from 
Table \ref{CeffTable}.

Eqn.(\ref{Sscat2mat}) {\em includes all multiple scattering interactions\/} within ${ \mit 
\Omega }_{j}$ to within the {\em weak Rayleigh approximation\/}. This technique is 
sometimes called {\em renormalization\/}. 
The technique is sometimes 
referred to as a {\em self avoiding renormalization\/} in quantum field 
theory, and is also implicit in the {\em redefinition of a scalar potential as 
an equivalent mass operator\/} in various applications of stochastic wave 
theory [e.g. {\em Sobczyk\/}~(1985)]. The self avoiding renormalization has important implications 
for the convergence of the {\em multiple scattering series\/}. This later 
point has been noted in {\em Schuster\/}~(1985), where an 
analogous strategy was employed to solve the {\em boundary integral 
equations\/} for a multiple inclusion problem. 



In practice the renormalization (\ref{dCeff0}) is 
effected using the matrix inversion
\begin{samepage}
\begin{eqnarray}  \label{dCeff}
\lefteqn{{\mit \Delta}{\bmi C}_{eff}({\bmi 
R}_j;k_y)=
}\\
&&
{\mit \Delta}{\bmi C}({\bmi 
R}_j;k_y)\cdot\left[ {\bmi I}_8 - \left(\lim\limits_{{\mit \Omega }_{j} \rightarrow   
0}\int_{{\mit \Omega }_{j}}{\tilde{\bmi  P}_{r,s}}
({\bmi R}_j,{\bmi  R}_j+{\bmi  r}_j){\rm d}{\bmi r}_j\right) \cdot{\mit \Delta}{\bmi C}({\bmi 
R}_j;k_y) \right]^{-1}\nonumber\\
&&
=\left[\begin{array}{cccccccc}
 -{\mit \Delta \rho}\,{\omega^2} &0&0&0&0&0&0&0\\
 0&{\mit \Delta}\mu\,{k_{y}^2} - {\mit \Delta} \rho\,{\omega^2} &0&0&0&0&0&
 -i\,{  k_y}\,{\mit \Delta \mu}\\
 0&0&{\mit \Delta}\tilde{\bmi C}_{33}&-{\mit \Delta}\tilde{\bmi C}_{43}&0&0&-{\mit \Delta}\tilde{\bmi C}_{73}&0\\
 0&0&{\mit \Delta}\tilde{\bmi C}_{43}&{\mit \Delta}\tilde{\bmi C}_{44}&0&0&{\mit \Delta}\tilde{\bmi C}_{74}&0\\
 0&0&0&0&{\mit \Delta}\tilde{\bmi C}_{55}&0&0&0\\
 0&0&0&0&0&{\mit \Delta}\tilde{\bmi C}_{66}&0&0\\
 0&0&{\mit \Delta}\tilde{\bmi C}_{73}&{\mit \Delta}\tilde{\bmi C}_{74}&0&0&{\mit \Delta}\tilde{\bmi C}_{77}&0\\
 0&i\,{  k_y}\,{\mit \Delta \mu}&0&0&0&0&0&{\mit \Delta \mu}
  \end{array}\right]\nonumber
\end{eqnarray}
\end{samepage}
with the components given in Table \ref{contrastTable} and 
Table \ref{CeffTable}. All of the elements of ${\mit \Delta}{\bmi C}_{eff}({\bmi 
R}_j;k_y)$ in (\ref{dCeff}) are evaluated at ${\bmi R}_j$. 


%%%%%%%%%%%%%%%%%%%%%%%%%%%%%%%%%%%%%%%%%%%%%%%%%%%%%%%%%%%%%%%%%%%%%%%%%%%%%%%%%%%%%%%%%%%
%%%%%%%%%%%%%%%%%%%%%%%%%%%%%%%%%%%%%%%%%%%%%%%%%%%%%%%%%%%%%%%%%%%%%%%%%%%%%%%%%%%%%%%%%%%
%%%%%%%%%%%%%%%%%%%%%%%%%%%%%%%%%%%%%%%%%%%%%%%%%%%%%%%%%%%%%%%%%%%%%%%%%%%%%%%%%%%%%%%%%%%%%%%
%%%%%%%%%%%%%%%%%%%%%%%%%%%%%%%%%%%%%%%%%%%%%%%%%%%%%%%%%%%%%%%%%%%%%%%%%%%%%%%%%%%%%%%%%%%%%%%%
%%%%%%%%%%%%%%%%%%%%%%%%%%%%%%%%%%%%%%%%%%%%%%%%%%%%%%%%%%%%%%%%%%%%%%%%%%%%%%%%%%%%%%%%%%%%%%%
\section{The matrix Lippmann-Schwinger equation}
In this final section, we will derive a matrix analogue of the augmented 
Lippmann-Schwinger type integral equation (\ref{LSTrep4}). 
To discretize (\ref{LSTrep4}) into a matrix equation, we replace the 
integral term with a matrix multiplication.  
 The matrix analogue translates directly into a computer code, provided the reference 
medium is homogeneous. The more useful case where the reference medium is 
a stratified half-space is considered at length in chapter \ref{stratG}.

First, we will arrange the effective coupling matrix ${\mit \Delta}{\bmi C}{}_{eff}({\bmi 
R}_j;k_y)$ of (\ref{dCeff}) into a larger matrix structure. This matrix 
is block-diagonal because the coupling operator is {\em local\/} in the spatial 
domain.  

The second extension concerns the plane wave expansion of the propagators.
We have already developed  matrix representations of the propagators 
appropriate to a lattice of source and receiver points, in 
\S\ref{latticeP}. However, these representations will require minor 
modification because the partial wave states used in the construction of 
the lattice propagators in \S\ref{latticeP} were $6-$component 
displacement/traction vectors. In this chapter, we have augmented those 
states by the  additional components $\partial_x u_x, \partial_x u_y$, 
for the purposes expounded above. Consequently, we will need to augment 
the plane wave lattice propagators of \S\ref{latticeP} in a similar way.
It is important to note that the number of states, and the manner in 
which they evolve over space remains unchanged. The only change is that 
each state carries some extra information in the two additional 
components. Consequently, it is straightforward to find the augmented 
propagators, and we simply state the extensions. 



\subsection{The lattice representation of the effective, weak-Rayleigh 
coupling matrix}

We recall the definition of the {\em secondary source field lattice\/} 
introduced in \S\ref{latticeP} and depicted in Fig.\ref{lattice}. This lattice was 
comprised of $L$ rows of lattice points, with $N_x$ lattice points per 
row, to give a total of $L \times N_x$ points.

We now define a {\em lattice expansion\/} of the {\em effective coupling array\/} 
${\mit \Delta}{\bmi C}{}_{eff}({\bmi 
R}_j;k_y)$ that was defined in (\ref{dCeff}). We recall from \S\ref{latticeP} that the 
lattice points are numbered in the same order as words to be read on the 
written page, i.e. top to bottom, left to right.
Once again, we may also  label the lattice 
points as if they were the elements of a $2-$D array. For this purpose we 
re-introduce the matrix row index $1\le \rho \le L$ and the matrix column 
index $1 \le \chi \le N_x$, so that the global index $i$ is related to 
the matrix indices via
\[
i=(\rho -1) N_x + \chi, \quad 1 \le i \le L N_x, 
\quad 1 \le \rho \le L, \quad 1 \le \chi \le N_x
\] 
The same {\em caveat\/} regarding the use of $\rho$ and $\chi$ to indicate the row and column 
indices respectively applies here, i.e. the  $\rho$ and $\chi$ in this 
section,  are not to confused with the material parameters given by the 
same symbols.

If we denote the coupling matrix for lattice points in 
$\mbox{row}_{\rho}$, as
${\mit \Delta}{\bmi C}{}^{\rho}_{eff}(k_y)$, this array may be defined explicitly as the $N_x$ element 
block diagonal array
\begin{eqnarray} \label{dCBCrho}
\lefteqn{{\mit \Delta}{\CB C}{}^{\rho}_{eff}(k_y)
={\mit \Delta}^{2}\times}\\
&& \mbox{diag}\left[{\begin{array}{cccc}
{\mit \Delta}{\bmi C}{}_{eff}({\bmi R}_{{N}_{x}( \rho 
  -1)+1};{k}_{y})& {\mit \Delta}{\bmi C}{}_{eff}({\bmi R}_{{N}_{x}( \rho  
 -1)+2};{k}_{y})&\cdots & {\mit \Delta}{\bmi C}{}_{eff}({\bmi R}_{ \rho {  
N}_{  x}};{k}_{y})\end{array}}\right]\nonumber
\end{eqnarray} 
where ${\mit \Delta}^{2}$ is the area of the small square subregion,
with the $8 \times 8$ diagonal block ${\mit \Delta}{\bmi C}{}_{eff}$ 
given by (\ref{dCeff}) at the respective lattice point. 

Finally, we assemble the block diagonal arrays ${\mit \Delta}{\bmi 
C}{}^{\rho}_{eff}(k_y)$ from (\ref{dCBCrho}) for each row of the lattice 
into the $L-$element, block diagonal array ${\mit \Delta}{\CB C}_{eff}(k_y)$ defined as
\[
\!\!\!\!\!\!\fbox{
\begin{minipage} {146.0mm}
\begin{equation} \label{dCBC}
{\mit \Delta}{\CB C}_{eff}(k_y)
=\mbox{diag}\left[{\begin{array}{cccc}
{\mit \Delta}{\CB C}{}^{1}_{eff}(k_y)& 
{\mit \Delta}{\CB C}{}^{1}_{eff}(k_y)&
\cdots & 
{\mit \Delta}{\CB C}{}^{1}_{eff}(k_y) 
\end{array}}\right]
\end{equation} 
\vspace{2 mm}
\end{minipage}
}
\]


\subsection{The augmented propagators }
In \S\ref{latticeP},  we wrote down explicit matrix representations for the 
propagators appropriate for a regular lattice expansion of some virtual 
or secondary source distribution. We assumed that this secondary source 
was represented by a $6-$vector denoted by ${\bmi f}_j$ at each of the $L 
\times N_x$ lattice points. However,  in \S\ref{secEffScat} of the present 
chapter, we have employed integration by parts to derive the effective 
coupling matrix ${\mit \Delta}{\bmi C}{}_{eff}({\bmi 
R}_j;k_y)$ that was defined in (\ref{dCeff}). This later operation 
entailed transferring the $\partial_x$ operators in ${{\mit \Delta} {\bmi A}}({\bmi x};{\partial 
}_{x},{k}_{y})$ of (\ref{weakRay}) onto 
the source and receiver coordinates of the propagator. Due to the {\em 
bilinear nature\/} of the propagator, the integration by parts was 
entirely equivalent to augmenting the propagator as well as the coupling 
matrix. The augmented form was very convenient for numerical 
applications, because the $6 \times 6 $ {\em differential\/} coupling operator 
${{\mit \Delta} {\bmi A}}({\bmi x};{\partial 
}_{x},{k}_{y})$ defined in (\ref{dAdef}) was replaced by the $8 \times 8$ 
{\em multiplicative\/} coupling matrix ${\mit \Delta}{\bmi C}{}_{eff}({\bmi 
R}_j;k_y)$ that was defined in (\ref{dCeff}). Mathematically, we are 
approximating a singular integro/differential operator equation, by a 
non-singular compact integral equation. 

We now present an augmented version of the plane wave expansion of the
{\em lattice propagator\/} from   (\ref{GGlattice}) of 
\S\ref{latticeP}.
 Certainly, we will not duplicate the manifold steps 
that lead us to the original representation (\ref{GGlattice}). Instead, 
we merely point out that  (\ref{GGlattice}) should be replaced by
\[
\!\!\!\!\!\!\fbox{
\begin{minipage} {146.0mm}
\begin{samepage}
\begin{eqnarray}\label{CBPaug}
\!\!\!\!\!\!\!\!\!
\rf{\CB P}_{r,s} &\equiv&
\rf{\CB P}\!{}_{r,s}^{z}\left[ \cdot \right]+ \rf{\CB  P}\!{}_{r,s}^{x}\nonumber\\
&&\!\!\!\!\!\!= -i\hat{\bmi \Phi}{}^{-1}\left[  {\bmi I}_L \otimes\left\{\left({\bmi I}_{8}\otimes 
{\bf DFT}({\bf k}_{x}) \right)\hat{\bmi  D}_{z}({\bf k}_{x},{k}_{y})\right\}
 {\CB Q}^{z}\right.\nonumber\\
&&\qquad\qquad\qquad\qquad\qquad
\times \left.
 {\bmi I}_L \otimes \left\{
\hat{\bmi  D}^{T}_{z}(-{\bf k}_{x},-{k}_{y})\left({\bmi I}_{8}\otimes
{\bf DFT}^{T}(-{\bf k}_{x})  \right)  \hat{\bmi \Phi}\left[ \cdot \right] \right\}\right]\nonumber\\
&&
 -{i}{\bmi I}_L \otimes\left( {\bmi I}_{N_x} \otimes
\left\{\left({\bmi I}_{8}\otimes \overline{\bmi F}(0)  \right)
 \hat{\bmi  D}_{x}({\bf k}_{z},{k}_{y})\right\}
 {\CB Q}{}^{x}_{\rho \leftarrow \rho}
\right.\nonumber\\
&&\qquad\qquad\qquad\qquad\qquad
\left.
\times{\bmi I}_{N_x} \otimes \left\{\hat{\bmi  D}_{x}^{T}(-{\bf k}_{z},-{k}_{y})
\,\left({\bmi I}_{8}\otimes \overline{\bmi F}^{T}(0)  \right)
\right\} \right)
\end{eqnarray}
\end{samepage}
\vspace{1 mm}
\end{minipage}
}
\]
where the $8 N_k \times 6 N_k$ array $\hat{\bmi  D}_{z}$  is an augmented version
 of the $6 N_k \times 6 N_k$ array 
${\bmi  D}_{z}$ (\ref{Dsuperdef}) and is given in {\sc Box } {\bf 4.5}.
Similarly the $8 \bar{N}_k \times 6 \bar{N}_k$ array $\hat{\bmi  D}_{x}$  is an augmented version
 of the $6 \bar{N}_k \times 6 \bar{N}_k$ array 
${\bmi  D}_{x}$ (\ref{Dsuperdef}) and is given in {\sc Box } {\bf 4.6}. 
The plane wave propagators ${\CB Q}^{z}$ and ${\CB Q}{}^{x}_{\rho \leftarrow \rho}$
were given by (\ref{CBQdef}) and (\ref{CBQxdef}) respectively.
The DFT matrix  ${\bf DFT}$  was given by (\ref{DFTdef}), and the 
projection matrix $\overline{\bmi F}(0)$ was given by (\ref{Fbardef}).
The augmented upsampling and downsampling transformations
$\hat{\bmi \Phi}\left[ \cdot \right]$ and $\hat{\bmi \Phi}{}^{-1}\left[ \cdot \right]$
respectively, are defined in Alg.4.1. By incorporating the upsampling and 
down sampling operators within the definition of $\rf{\CB 
P}\!\!{}_{r,s}^{z}\left[ \cdot \right]$ in (\ref{CBPaug}),  $\rf{\CB 
P}\!\!{}_{r,s}$ operates directly on $8 L N_x-$vectors.

%%%%%%%%%%%%%%%%%%%%%%%%%%%%%%%%%%%%%%%%%%%%%%%%%%%%%%%%%%%%%%%%%%%%%%%%%%%%%%%%%%%%%%%%%%%%%%%
%%%%%%%%%%%%%%%%%%%%%%%%%%%%%%%%%%%%%%%%%%%%%%%%%%%%%%%%%%%%%%%%%%%%%%%%%%%%%%%%%%%%%
%%%%%%%%%%%%%%%%%%%%%%%%%%%%%%%%%%%%%%%%%%%%%%%%%%%%%%%%%%%%%%%%%%%%%%%%%%%%%%%%%%%%%%%%%%%
\begin{figure}
\fbox{
\begin{minipage} {140.0mm}
{\bf  Algorithm}~{\bf 4.1} {\sc The augmented analogues $\hat{\bmi \Phi}[\cdot]$ and 
$\hat{\bmi \Phi}{}^{-1}[\cdot]$ of the sampling operators from {\bf  
Algorithm}~{\bf 3.1} }
\vspace{3mm}

\begin{tt}

Complex $\tilde{\CB F}_{\rho}(N_k,8), {\CB F}_{\rho}(8,N_x)$\newline
C0=(0,0)


\begin{itemize}
\item Initialize  $\tilde{\CB F}_{\rho}$


\begin{itemize}
\item For j=1,8
\begin{itemize}
\item For k=1,$N_k$
\begin{itemize}
\item $\tilde{\CB F}_{\rho}$(k,j)=C0
\end{itemize}
\item End !k-loop
\end{itemize}
\item End !j-loop
\end{itemize}

\vspace{1 mm}

\item The upsampling  mapping $\tilde{\CB F}_{\rho}=\hat{\bmi \Phi}[{\CB F}_{\rho}]$

\begin{itemize}
\item For j=1,8
\begin{itemize}
\item For k=1,$N_x$
\begin{itemize}
\item $\tilde{\CB F}_{\rho}$((k-1)$N_{\mit 
\Delta}+$1,j)=${\CB F}_{\rho}$(j,k)
\end{itemize}
\item End !k-loop
\end{itemize}
\item End !j-loop
\end{itemize}



\end{itemize}

\vspace{2 mm}

\begin{itemize}
\item The downsampling  mapping ${\CB F}_{\rho}=\hat{\bmi \Phi}{}^{-1}[\tilde{\CB F}_{\rho}]$
\begin{itemize}
\item For j=1,8
\begin{itemize}
\item For k=1,$N_x$
\begin{itemize}
\item ${\CB F}_{\rho}$(j,k)$=\tilde{\CB F}_{\rho}$((k-1)$N_{\mit 
\Delta}+$1,j)
\end{itemize}
\item End !k-loop
\end{itemize}
\item End !j-loop
\end{itemize}
\end{itemize}

\end{tt} 

\vspace{2 mm}
\end{minipage}
}
\end{figure}
%%%%%%%%%%%%%%%%%%%%%%%%%%%%%%%%%%%%%%%%%%%%%%%%%%%%%%%%%%%%%%%%%%%%%%%%%%%%%%%%%%%%%%%%%%%
%%%%%%%%%%%%%%%%%%%%%%%%%%%%%%%%%%%%%%%%%%%%%%%%%%%%%%%%%%%%%%%%%%%%%%%%%%%%%%%%%%%%%

\begin{figure}
\fbox{
\begin{minipage} {140.0mm}
{\sc  Box} ${\bf  4.5}$ {\sc The augmented up/down plane wave states\/}\newline

\begin{samepage}
\begin{eqnarray} \label{hatDsuperdef}
\lefteqn{\hat{\bmi D}_{z}({\bf k}_{x},k_{y})=}\\
&&\!\!\!\!\!\!\!\!\!\!\!\!
\left[\begin{array}{cccccc}
\left|+,P, {\bf k}_{x},k_{y}  \right)&\left|+,S, {\bf k}_{x},k_{y}  
\right)&\left|+,H, {\bf k}_{x},k_{y}  \right)&
\left|-,P, {\bf k}_{x},k_{y}  \right)&\left|-,S, {\bf k}_{x},k_{y}  
\right)&\left|-,H, {\bf k}_{x},k_{y}  \right)
\end{array}\right]\nonumber
\end{eqnarray}
\end{samepage}
is defined in terms of the partial wave states
\begin{itemize}
\item the $\left|\pm,P, {\bf k}_x,k_y  \right)$ partial wave-vectors 
are the   ${\stackrel{\mbox{\small downgoing}}{{}_{\mbox{\small 
upgoing}}}}$  $P-$waves 
\begin{equation} \label{hatPeigen}
\left|\pm,P, {\bf k}_x ,k_y \right)=
\sqrt{\frac{{ \Delta} k_{x}}{{4 \pi {\rf{\rho}} {\omega 
}^{  2}{\bf  k}_{  
z, c }}}}\,
\left[\begin{array}{c} \pm i\,{\bf k}_{z,P}\\i\,{\bf k}_{x}\\
i\,k_{y}\,{\bmi I}\\
  {\rf{\rho}} \,\left( 2\,{{\rf{\beta}}{}^{2}}\,{{\bf k}_{x}^2} + ( 2\,{{\rf{\beta}}{}^{2}}\,{k_{y}^2} - 
     {\omega^2}){\bmi I} \right) \\ \mp 2\,{\rf{\rho}}\,{{\rf{\beta}}{}^{2}}\,{\bf k}_{x}\,{\bf k}_{z,P} \\
  \mp 2\,{\rf{\rho}}\,{{\rf{\beta}}{}^{2}}\,k_{y}\,{\bf k}_{z,P}\\-{{\bf k}_{x}^2}\\-{k_{y}}\,{{\bf k}_{x}} \end{array}\right]
\end{equation}

\item the $\left|\pm,S, {\bf k}_x ,k_y \right)$ partial wave-vector 
are the   ${\stackrel{\mbox{\small downgoing}}{{}_{\mbox{\small 
upgoing}}}}$  quasi-$SV-$waves 
\begin{equation} \label{hatSVeigen}
\left|\pm,S, {\bf k}_x ,k_y \right)=\frac{\omega }{\rf{\beta} \hat{K}_{S}}\,
\sqrt{\frac{{ \Delta} k_{x}}{{4 \pi {\rf{\rho}} {\omega 
}^{  2}{\bf  k}_{  
z, c }}}}\,
\left[\begin{array}{c} 
i\,{\bf k}_{x}\\
\mp i\,{\bf k}_{z,S}\\ {\bf 0} \\
\mp 2\,{\rf{\rho}}\,{{\rf{\beta}}{}^{2}}\,{\bf k}_{x}\,{\bf k}_{z,S} \\
  {\rf{\rho}} \,\left(({\omega^2} - {{\rf{\beta}}{}^{2}}\,{k_{y}^2}) {\bmi I}-2\,{{\rf{\beta}}{}^{2}}\,{{\bf k}_{x}^2}   \right) \\
  -{\rf{\rho}}\,{{\rf{\beta}}{}^{2}}\,k_{y}\,{\bf k}_{x}\\ \pm {{\bf k}_{x}}\,{\bf 
  k}_{z,S}\\{\bf 0} \end{array}\right]
\end{equation}

\item the $\left|\pm,H, {\bf k}_x ,k_y \right)$ partial wave-vectors 
are the   ${\stackrel{\mbox{\small downgoing}}{{}_{\mbox{\small 
upgoing}}}}$  quasi-$SH-$waves 
\begin{equation} \label{hatSHeigen}
\left|\pm,H, {\bf k}_x ,k_y \right)=\frac{1 }{\hat{K}_{H}}\,
\sqrt{\frac{{ \Delta} k_{x}}{{4 \pi {\rf{\rho}} {\omega 
}^{  2}{\bf  k}_{  
z, c }}}}\,
\left[\begin{array}{c} 
\mp k_{y}\,{\bf k}_{z,H} \\-k_{y}\,{\bf k}_{x} \\
  \hat{K}_{H}^2 {\bmi I}\\
-2\,i\,{\rf{\rho}} \,{{\rf{\beta}}{}^{2}}\,k_{y}\,{{\bf k}_{z,H}^2}\\
  \mp 2\,i\,{\rf{\rho}}\,{{\rf{\beta}}{}^{2}}\,k_{y}\,{\bf k}_{x}\,{\bf k}_{z,H} \\
  \pm i\,{\rf{\rho}}\,\left(  {\omega^2} -2\,{{\rf{\beta}}{}^{2}}\,{k_{y}^2} 
  \right)\,{\bf k}_{z,H} \\
  -i\,{k}_{y}\,{\bf k}^{2}_{x}\\i\,\hat{K}{}^{2}_{H}\,{\bf k}_{x} \end{array}\right]
\end{equation}
\end{itemize}
where  ${\bf  k}_{x}$ and ${\bf  k}_{z,c}$ were given in
(\ref{gKx}) and (\ref{gKz}) respectively. 
\vspace{2 mm}
\end{minipage}
}
\end{figure}



\begin{figure}
\fbox{
\begin{minipage} {140.0mm}
{\sc  Box} ${\bf  4.6}$ {\sc The augmented left/right plane wave states\/}\newline

\begin{eqnarray} \label{hatDxsuperdef}
\lefteqn{\hat{\bmi D}_{x}({\bf k}_{z},k_{y})=}\\
&&\!\!\!\!\!\!\!\!\!\!\!\!
\left[\begin{array}{cccccc}
\overline{\left|+,P, {\bf k}_{z},k_{y}  \right)}&\overline{\left|+,S, {\bf k}_{z},k_{y}  
\right)}&\overline{\left|+,H, {\bf k}_{z},k_{y}  \right)}&
\overline{\left|-,P, {\bf k}_{z},k_{y}  \right)}&\overline{\left|-,S, {\bf k}_{z},k_{y}  
\right)}&\overline{\left|-,H, {\bf k}_{z},k_{y}  \right)}
\end{array}\right]\nonumber
\end{eqnarray}
is defined in terms of the partial wave states
\begin{itemize}
\item the $8 \overline{N}_k \times \overline{N}_k$ $P$-wave partition
\begin{equation} \label{hatPxeigenSuper}
\overline{\left|\pm,P, {\bf k}_z,k_{y}  \right)}=
\sqrt{\frac{{ \Delta} k_{z}}{{4 \pi {\rf{\rho}} {\omega 
}^{  2}{\bf  k}_{  
x, P }}}}
\left[\begin{array}{c} i\,{\bf k}_{z}\\ \pm i\,{\bf k}_{x,P}\\
i\,k_{y}\,{\bmi I}\\
  {\rf{\rho}} \,\left( 2\,{{\rf{\beta}}{}^{2}}\,{{\bf k}_{x,P}^2} + ( 2\,{{\rf{\beta}}{}^{2}}\,{k_{y}^2} - 
     {\omega^2}){\bmi I} \right) \\ \mp 2\,{\rf{\rho}}\,{{\rf{\beta}}{}^{2}}\,{\bf k}_{x,P}\,{\bf k}_{z} \\
  - 2\,{\rf{\rho}}\,{{\rf{\beta}}{}^{2}}\,k_{y}\,{\bf k}_{z}
  \\ -\,{\bf k}^{2}_{x,P}\\
\mp \,k_{y}\,{\bf k}_{x,P} \end{array}\right]
\end{equation}

\item  the $8 \overline{N}_k \times \overline{N}_k$ $SV$-wave partition
\begin{equation} \label{hatSVxeigenSuper}
\overline{\left|\pm,S, {\bf k}_z,k_{y}  \right)}=
\frac{\omega }{\rf{\beta} \hat{K}_{S}}\,\sqrt{\frac{{ \Delta} k_{z}}{{4 \pi {\rf{\rho}} {\omega 
}^{  2}{\bf  k}_{  
x, S }}}}\,\left[\begin{array}{c} 
\pm i\,{\bf k}_{x,S}\\
- i\,{\bf k}_{z}\\ {\bf 0} \\
\mp 2\,{\rf{\rho}}\,{{\rf{\beta}}{}^{2}}\,{\bf k}_{x,S}\,{\bf k}_{z} \\
  {\rf{\rho}} \,\left(({\omega^2} - {{\rf{\beta}}{}^{2}}\,{k_{y}^2}) {\bmi 
  I}-2\,{{\rf{\beta}}{}^{2}}\,{{\bf k}_{x,S}^2}   \right) \\
 \mp {\rf{\rho}}\,{{\rf{\beta}}{}^{2}}\,k_{y}\,{\bf k}_{x,S}
 \\
\pm \,{\bf k}_{z} {\bf k}_{x,S}\\ {\bf 0} \end{array}\right]
\end{equation}

\item the $8 \overline{N}_k \times \overline{N}_k$ $SH$-wave partition
\begin{equation} \label{hatSHxeigenSuper}
\overline{\left|\pm,H, {\bf k}_z, k_{y}  \right)}=\frac{1 }{\hat{K}_{H}}\,
\sqrt{\frac{{ \Delta} k_{z}}{{4 \pi {\rf{\rho}} {\omega 
}^{  2}{\bf  k}_{  
x, H }}}}\,\left[\begin{array}{c} 
- k_{y}\,{\bf k}_{z} \\ \mp k_{y}\,{\bf k}_{x,H} \\
  \hat{K}_{H}^2 {\bmi I}\\
-2\,i\,{\rf{\rho}} \,{{\rf{\beta}}{}^{2}}\,k_{y}\,{{\bf k}_{z}^2}\\
  \mp 2\,i\,{\rf{\rho}}\,{{\rf{\beta}}{}^{2}}\,k_{y}\,{\bf k}_{x,H}\,{\bf k}_{z} \\
   i\,{\rf{\rho}}\,\left(  {\omega^2} -2\,{{\rf{\beta}}{}^{2}}\,{k_{y}^2} 
  \right)\,{\bf k}_{z}
  \\ -i\, k_{y}\,{\bf k}^{2}_{x,H} \\
  \pm i\, \hat{K}_{H}^2 {\bf k}_{x,H}  \end{array}\right]
\end{equation}

\end{itemize}
where  ${\bf  k}_{x}$ and ${\bf  k}_{z,c}$ were given in
(\ref{gbarKx}) and (\ref{gbarKz}) respectively.
\vspace{2 mm}
\end{minipage}
}
\end{figure}


%%%%%%%%%%%%%%%%%%%%%%%%%%%%%%%%%%%%%%%%%%%%%%%%%%%%%%%%%%%%%%%%%%%%%%%%%%%%%%%%%%%%%%%%%%%
%%%%%%%%%%%%%%%%%%%%%%%%%%%%%%%%%%%%%%%%%%%%%%%%%%%%%%%%%%%%%%%%%%%%%%%%%%%%%%%%%%%%%%%%%%%
%%%%%%%%%%%%%%%%%%%%%%%%%%%%%%%%%%%%%%%%%%%%%%%%%%%%%%%%%%%%%%%%%%%%%%%%%%%%%%%%%%%%%%%%%%%%%%%
%%%%%%%%%%%%%%%%%%%%%%%%%%%%%%%%%%%%%%%%%%%%%%%%%%%%%%%%%%%%%%%%%%%%%%%%%%%%%%%%%%%%%%%%%%%%%%%%
%%%%%%%%%%%%%%%%%%%%%%%%%%%%%%%%%%%%%%%%%%%%%%%%%%%%%%%%%%%%%%%%%%%%%%%%%%%%%%%%%%%%%%%%%%%%%%%
%%%%%%%%%%%%%%%%%%%%%%%%%%%%%%%%%%%%%%%%%%%%%%%%%%%%%%%%%%%%%%%%%%%%%%%%%%%%%%%%%%%%%
%%%%%%%%%%%%%%%%%%%%%%%%%%%%%%%%%%%%%%%%%%%%%%%%%%%%%%%%%%%%%%%%%%%%%%%%%%%%%%%%%%%%%%%%%%
%%%%%%%%%%%%%%%%%%%%%%%%%%%%%%%%%%%%%%%%%%%%%%%%%%%%%%%%%%%%%%%%%%%%%%%%%%%%%%%%%%%%%%%%%%%
%%%%%%%%%%%%%%%%%%%%%%%%%%%%%%%%%%%%%%%%%%%%%%%%%%%%%%%%%%%%%%%%%%%%%%%%%%%%%%%%%%%%%%%%%%%
%%%%%%%%%%%%%%%%%%%%%%%%%%%%%%%%%%%%%%%%%%%%%%%%%%%%%%%%%%%%%%%%%%%%%%%%%%%%%%%%%%%%%%%%%%%%%%%
%%%%%%%%%%%%%%%%%%%%%%%%%%%%%%%%%%%%%%%%%%%%%%%%%%%%%%%%%%%%%%%%%%%%%%%%%%%%%%%%%%%%%%%%%%%%%%%%

\subsection{The matrix formulation of the governing equation }
Using the propagator ${\rf{\CB P}}_{r,s}$ from (\ref{CBPaug}) and the coupling 
matrix ${\mit \Delta} {\CB C}_{eff}$ from (\ref{dCBC}), 
the matrix analogue of the governing equation (\ref{LSTrep4}) may be written as
\begin{equation} \label{LSTrep5}
{{\CB G}}_{r}={\rf{\CB G}}_{r}+{\rf{\CB P}}_{r,s}{\mit \Delta} {\CB C}_{eff} {\CB G}_{r}
\end{equation} 
where ${\rf{\CB G}}_{r}$ is the augmented reference Green's tensor whose 
explicit form depends on the location and nature of the seismic source.



